\section{Supplementary Information}\label{sec:supp}

% =========================================================
% Background (kept, uncited but pedagogical)
% =========================================================

\subsection{Background on \texorpdfstring{$\mathrm{Cl}(3,0)$}{Cl(3,0)} and the geometric product}\label{supp:prelim}

The real Clifford algebra $\mathrm{Cl}(3,0)$ is the eight-dimensional associative unital algebra generated by an orthonormal basis $\{e_1, e_2, e_3\}$ of $\mathbb{R}^3$ subject to $e_ie_j + e_je_i = 2\delta_{ij}$ \citep{lounesto2001clifford,doran2003geometric}. A basis of $\mathrm{Cl}(3,0)$ ordered by grade is $\{1,\; e_1, e_2, e_3,\; e_{12}, e_{13}, e_{23},\; e_{123}\}$, corresponding respectively to a scalar (grade 0), three vectors (grade 1), three bivectors (grade 2), and a pseudoscalar (grade 3). Any multivector $X \in \mathrm{Cl}(3,0)$ decomposes uniquely as $X = X^{(0)} + X^{(1)} + X^{(2)} + X^{(3)}$ with $X^{(k)}$ the grade-$k$ component.

The geometric product is the unique bilinear, associative, unital product satisfying the generator relations above. For two vectors $u, v \in \mathbb{R}^3$ embedded in $\mathrm{Cl}(3,0)$ it decomposes as $uv = u \cdot v + u \wedge v$, where $u \cdot v = \tfrac{1}{2}(uv + vu) = \sum_i u_iv_i$ is the symmetric scalar part and $u \wedge v = \tfrac{1}{2}(uv - vu) = \sum_{i<j}(u_iv_j - u_jv_i)e_{ij}$ is the antisymmetric bivector part.

The pseudoscalar $I = e_1e_2e_3$ satisfies $I^2 = -1$, so $I^{-1} = -I$. The Hodge dual $\star: \Lambda^k\mathbb{R}^3 \to \Lambda^{3-k}\mathbb{R}^3$ can be defined as $\star\alpha = \alpha\, I^{-1}$. Direct computation yields the sign convention
\[
\star e_{12} = +e_3,\qquad \star e_{13} = -e_2,\qquad \star e_{23} = +e_1,
\]
which is the standard right-handed orientation of \citet{doran2003geometric,lounesto2001clifford} under which $\star(u\wedge v)$ equals the ordinary cross product $u\times v$. Under $\mathrm{SO}(3)$, $\star$ is equivariant. Under $\mathrm{O}(3)$, the Hodge dual of a bivector is an axial (pseudo)vector: bivectors are invariant under spatial inversion ($e_ie_j \to (-e_i)(-e_j) = e_ie_j$), so the Hodge vector inherits parity $+1$ while a polar vector carries parity $-1$; see Supplementary~\ref{supp:so3_o3}.

The geometric product is $\mathrm{O}(3)$-equivariant under the rotor action: for any rotor $R = \exp(-\theta B/2)$ built from a unit bivector $B$ and any multivectors $a, b$, conjugation $a \mapsto RaR^{-1}$ preserves the product, $R(ab)R^{-1} = (RaR^{-1})(RbR^{-1})$ \citep{doran2003geometric}. This is the single inductive bias the Clifford track exploits.

% =========================================================
% Related-work cited sections
% =========================================================

\subsection{Architectural comparison with ICTP}\label{supp:ictp}

CliffordSTF and ICTP~\citep{zaverkin2024ictp} both use symmetric traceless Cartesian tensors as the primary carrier of $L \geq 2$ information, and both prove equivariance and tracelessness of the resulting layers. The two architectures nevertheless differ in four substantive ways that correspond to distinct design philosophies.

\textbf{Algebraic substrate.} ICTP is built entirely from irreducible Cartesian tensors and their products; spherical tensors do not appear. CliffordSTF runs a Clifford $\mathrm{Cl}(3,0)$ multivector track in parallel with the STF tracks, and the Clifford track carries both scalars and pseudoscalars that have a natural parity structure through the grade involution, as well as bivectors that expose an extra $L{=}1$ axial channel via the Hodge star. The STF tracks in CliffordSTF supply the $L{=}2$ and $L{=}3$ components that the Clifford geometric product cannot reach, not the entire equivariant substrate.

\textbf{Bilinear operation.} ICTP uses $\nu$-fold tensor products of irreducible Cartesian tensors as its core bilinear, following MACE-style many-body expansion. CliffordSTF uses two different bilinears in tandem: the Clifford geometric product acting on the 8-dimensional multivector via an $8 \times 8$ Cayley table, and a small set of closed-form symmetric-traceless contractions acting on the STF tracks. The Clifford product handles scalar, vector, bivector, and pseudoscalar channels with a single contraction; the STF operations handle rank-2 and rank-3 tensors through explicit symmetrisation and trace removal. The two are coupled bijectively to Clebsch--Gordan channels as specified in Theorem~\ref{thm:cg_completeness} and its higher-order extensions.

\textbf{Cross-track coupling.} In ICTP, interactions between rank-$L_1$ and rank-$L_2$ tensors happen through a single irreducible Cartesian tensor product operation. In CliffordSTF the Clifford and STF tracks are coupled bidirectionally: the STF$_2$ track is generated from grade-1 vector features via $\mathrm{STF}_2(\hat{r}_{ij}, h_j^{\mathrm{vec}})$, and is fed back into the grade-1 vector component of the Clifford message through the $S \cdot v$ contraction. Ablating this coupling (the \texttt{stf2\_no\_cross} variant) causes a substantial performance regression compared to the base Clifford model, demonstrating that the STF track by itself does not carry a usable training signal without coupling back to the Clifford substrate.

\textbf{Equivariance scope.} ICTP proves $\mathrm{O}(3)$ equivariance (rotations and reflections), with parity tracked explicitly through the product structure. CliffordSTF proves $\mathrm{SO}(3)$ equivariance; under $\mathrm{O}(3)$, the polar-vector grade-1 component and the axial-vector Hodge-of-bivector component would need distinct parity labels $1^-$ and $1^+$ respectively (Remark~\ref{rem:so3_vs_o3}, Supplementary~\ref{supp:so3_o3}). The current implementation concatenates the two in the force readout under the $\mathrm{SO}(3)$ framework, which is consistent under rotations but would need refactoring for full $\mathrm{O}(3)$.

These differences explain why the two architectures occupy different operating points. ICTP targets state-of-the-art accuracy on small-molecule benchmarks at the cost of a denser bilinear. CliffordSTF targets a minimal algebraic substrate with closed-form coupling, trading raw small-molecule accuracy for parameter efficiency on catalysis (OC20, OC22) and for a direct mechanistic handle on which design choices drive which gains, as demonstrated by our eleven-variant ablation (\S\ref{sec:results_scalar}).

\subsection{Reconciliation with prior Clifford-network grade structure}\label{supp:gatr}

\citet{brehmer2023geometric} builds a geometric algebra transformer in the four-dimensional projective algebra $\mathrm{Cl}(3,0,1)$. Under the $\mathrm{SO}(3)$ subgroup of its symmetry group, GATr's hidden multivectors decompose as two copies of $L{=}0$ (a scalar and a trivector) and two copies of $L{=}1$ (a vector and a bivector, axial under Hodge duality); no $L \geq 2$ irrep is represented inside the multivector. The same decomposition holds for the non-projective $\mathrm{Cl}(3,0)$ used here, because the extra null basis element of the projective algebra carries only translational information and does not contribute $L \geq 2$ content. Our diagnostic is therefore about the $\mathrm{SO}(3)$-irrep content of the multivector substrate \emph{at the linear bilinear level}, not about the choice of projective versus non-projective algebra: both carry the same $L{=}2$ deficit. \citet{ruhe2023cgenn} shows that $L \geq 2$ content can be reconstructed \emph{nonlinearly} through repeated geometric products, but this indirect access incurs the depth cost and gradient pathologies that our direct STF tracks avoid.

% =========================================================
% Method-cited sections (in citation order)
% =========================================================

\subsection{\texorpdfstring{$\mathrm{SO}(3)$}{SO(3)} versus \texorpdfstring{$\mathrm{O}(3)$}{O(3)} equivariance}\label{supp:so3_o3}

All guarantees stated in \S\ref{sec:method} are $\mathrm{SO}(3)$-equivariance guarantees. Under $\mathrm{SO}(3)$, the $L{=}1$ representation is unique up to isomorphism, so polar vectors (grade 1 of $\mathrm{Cl}(3,0)$) and axial vectors (Hodge duals of grade-2 bivectors) transform identically and may be freely concatenated in the force readout. Under $\mathrm{O}(3)$ they differ by a factor of $\det(R)$ under the action of $R \in \mathrm{O}(3)\setminus\mathrm{SO}(3)$, corresponding to the parity labelling summarised in Table~\ref{tab:parity}; the convention follows Jackson's \emph{Classical Electrodynamics} \S6.10 and the \texttt{e3nn} library convention \citep{geiger2022e3nn}.

\begin{table}[h]
\centering
\caption{$\mathrm{O}(3)$ parity labels for the tracks used in CliffordSTF. The general rule is $\mathrm{STF}_\ell$ built from polar vectors has parity $(-1)^\ell$, matching the parity of the spherical harmonic $Y_{\ell m}(-\hat{n}) = (-1)^\ell Y_{\ell m}(\hat{n})$.}\label{tab:parity}
\small
\begin{tabular}{lccc}
\toprule
Object & Physics analog & Parity & $\mathrm{O}(3)$ label \\
\midrule
Grade-0 scalar & charge $q$ & $+$ & $0^+$ \\
Grade-1 vector $e_i$ & position $\mathbf{r}$ & $-$ & $1^-$ (polar) \\
Grade-2 bivector $e_{ij}$ / Hodge dual & angular momentum $\mathbf{L}{=}\mathbf{r}\times\mathbf{p}$ & $+$ & $1^+$ (axial) \\
Grade-3 pseudoscalar $I$ & triple product $\mathbf{a}\cdot(\mathbf{b}\times\mathbf{c})$ & $-$ & $0^-$ \\
$\mathrm{STF}_2$ from polar vectors & quadrupole $Q_{ij}$ & $+$ & $2^+$ \\
$\mathrm{STF}_3$ from polar vectors & octupole $O_{ijk}$ & $-$ & $3^-$ \\
\bottomrule
\end{tabular}
\end{table}

Extending to full $\mathrm{O}(3)$-equivariance would require treating the grade-1 polar-vector feature and the Hodge-of-bivector axial-vector feature as distinct $1^-$ and $1^+$ tracks rather than concatenating them, and likewise distinguishing $2^+$ and $3^-$ tracks. All bilinear interactions in the architecture respect this parity grading when made explicit; the current implementation simply chooses not to distinguish $1^+$ from $1^-$, which is consistent under $\mathrm{SO}(3)$ but would conflate them under $\mathrm{O}(3)$.

\subsection{Derivation of the STF coefficients}\label{supp:stf_coefficients}

We derive the three rational coefficients appearing in the STF$_2$ and STF$_3$ formulas of \S\ref{sec:stf_channels}.

\textbf{STF$_2$ detracing coefficient $\tfrac{1}{3}$.} In $d$ dimensions, the symmetric traceless part of a rank-2 tensor requires subtracting $\tfrac{1}{d} \delta_{ij} T_{kk}$ to enforce $\sum_i T_{ii}^{\text{sym-traceless}} = 0$. Direct check in $d=3$: $\delta_{ii} \mathrm{STF}_2(u,v)_{ii} = (u\cdot v) - \tfrac{1}{3}(3)(u\cdot v) = 0$.

\textbf{STF$_3$ trace vector coefficient $\tfrac{2}{3}$.} Since $S_{ij}$ is symmetric traceless ($S_{ii}=0$), the trace of $T^{\text{sym}}_{ijk}$ over any pair of indices reduces to two of the three summands:
\[
T^{\text{sym}}_{iik} = \tfrac{1}{3}(S_{ii}v_k + S_{ik}v_i + S_{ik}v_i) = \tfrac{1}{3}(0 + 2\,S_{kl}v_l) = \tfrac{2}{3}\,S_{kl}v_l.
\]
By full symmetry, $T^{\text{sym}}_{iki} = T^{\text{sym}}_{kii} = A_k$ as well.

\textbf{STF$_3$ detracing coefficient $\tfrac{1}{5}$.} Set $U_{ijk} = T^{\text{sym}}_{ijk} - \alpha(\delta_{ij}A_k + \delta_{ik}A_j + \delta_{jk}A_i)$ and require $U_{iik} = 0$. Computing the trace:
\[
U_{iik} = A_k - \alpha(3A_k + A_k + A_k) = A_k - 5\alpha A_k \implies \alpha = \tfrac{1}{5}.
\]
More generally, in $d$ dimensions the per-delta first-trace coefficient at rank $\ell$ is $1/(2\ell + d - 4)$, which reduces to $1/(2\ell - 1)$ in $d=3$ and yields $\tfrac{1}{3}$ at $\ell{=}2$, $\tfrac{1}{5}$ at $\ell{=}3$, $\tfrac{1}{7}$ at $\ell{=}4$ \citep{applequist1989traceless,thorne1980multipole}. The fully detraced closed form for arbitrary rank is \citet[Theorem~2]{applequist1989traceless}.

\textbf{Degrees-of-freedom check for STF$_3$.} The symmetric rank-3 tensor $T^{\text{sym}}$ in $\mathbb{R}^3$ has $\binom{3+3-1}{3} = 10$ independent components; subtracting the 3-component trace vector $A$ leaves $10 - 3 = 7 = 2\ell + 1$ for $\ell = 3$, matching the dimension of the $L{=}3$ irrep of $\mathrm{SO}(3)$.

\subsection{Explicit coordinate formulas for the constructive CG maps}\label{supp:cg_proof}

We enumerate every constructive map used in the architecture and verified in Supplementary~\ref{supp:cg_test}. Symmetric traceless projection is a textbook construction \citep{coope1965irreducible,applequist1989traceless,stone2013theory}; rank-3 STF projectors appear in \citet{thorne1980multipole}.

\textbf{$\mathbf{1} \otimes \mathbf{1} \to \mathbf{0} \oplus \mathbf{1} \oplus \mathbf{2}$.}
For $u, v \in \mathbb{R}^3$:
\begin{itemize}
    \item $L{=}0$: $(u \otimes v)_{L=0} = u \cdot v = u_iv_i$.
    \item $L{=}1$: $(u \otimes v)_{L=1} = \star(u \wedge v)_i = \varepsilon_{ijk}u_jv_k$, equivalently the cross product $u \times v$.
    \item $L{=}2$: $(u \otimes v)_{L=2} = \mathrm{STF}_2(u,v)_{ij} = \tfrac{1}{2}(u_iv_j + u_jv_i) - \tfrac{1}{3}(u\cdot v)\delta_{ij}$.
\end{itemize}

\textbf{$\mathbf{2} \otimes \mathbf{1} \to \mathbf{1} \oplus \mathbf{2} \oplus \mathbf{3}$.}
For a symmetric traceless $S \in \mathbb{R}^{3\times 3}$ and $v \in \mathbb{R}^3$:
\begin{itemize}
    \item $L{=}1$: $(S \cdot v)_i = S_{ij}v_j$.
    \item $L{=}2$: symmetric traceless part of $\varepsilon_{ikl}S_{jl}v_k$, i.e.\ $[\mathrm{sym}(S\times v) - \tfrac{1}{3}\mathrm{tr}(\cdot)\mathbb{I}]$; explicit 5-component formulas are in the code release (\texttt{contract\_stf2\_vec\_to\_stf2}).
    \item $L{=}3$: $\mathrm{STF}_3(S, v)_{ijk} = S_{(ij}v_{k)} - \tfrac{1}{5}(\delta_{ij}A_k + \delta_{ik}A_j + \delta_{jk}A_i)$ with $A_k = \tfrac{2}{3}S_{kl}v_l$; the trace coefficient $\tfrac{2}{3}$ follows from tracelessness of $S$ and the detracing coefficient $\tfrac{1}{5}$ from $1/(2\ell{-}1)|_{\ell=3}$ in $d{=}3$ \citep{applequist1989traceless,thorne1980multipole}.
\end{itemize}

\textbf{$\mathbf{2} \otimes \mathbf{2} \to \mathbf{0} \oplus \mathbf{1} \oplus \mathbf{2} \oplus \mathbf{3}$ (with $L{=}4$ omitted).}
For symmetric traceless $S_1, S_2$:
\begin{itemize}
    \item $L{=}0$: $\langle S_1, S_2 \rangle_F = S_{1,ij}S_{2,ij}$ (Frobenius inner product).
    \item $L{=}1$: antisymmetric contraction $\varepsilon_{ijk}S_{1,jl}S_{2,kl}$.
    \item $L{=}2$: symmetric traceless projection of $S_{1,ik}S_{2,kj}$.
    \item $L{=}3$: an analogous symmetric-then-detraced construction treating one operand as a vector-indexed rank-2 tensor.
    \item $L{=}4$: outside our $L_{\max}{=}3$ cut and not used.
\end{itemize}

\textbf{$\mathbf{3} \otimes \mathbf{1} \to \mathbf{2} \oplus \mathbf{3}$ (with $L{=}4$ omitted).}
For a symmetric traceless rank-3 tensor $T$ and vector $v$:
\begin{itemize}
    \item $L{=}2$: $(T \cdot v)_{ij} = T_{ijk}v_k$ followed by trace removal (\texttt{contract\_stf3\_vec\_to\_stf2}).
    \item $L{=}3$: analogous symmetric-then-detraced construction.
    \item $L{=}4$: omitted.
\end{itemize}

Each map is a polynomial in its inputs with constant rational coefficients. There are no learned Clebsch--Gordan coefficients, no Wigner-$D$ tables, and no calls to \texttt{e3nn} \citep{geiger2022e3nn}. All tensor products of $\mathrm{SO}(3)$ irreps above are multiplicity-free, consistent with the general $\mathrm{SU}(2)$ fact that $L_1 \otimes L_2 = \bigoplus_{L=|L_1-L_2|}^{L_1+L_2} L$ each with multiplicity one \citep{varshalovich1988quantum}.

\subsection{Numerical verification of the constructive CG maps and end-to-end equivariance}\label{supp:cg_test}

\textbf{Primitive-level verification.} We verify Theorem~\ref{thm:cg_completeness} and its higher-order extensions numerically by comparing each constructive map against a reference implementation based on Wigner $3j$-symbols acting on spherical-harmonic representations. For each admissible triple $(\ell_1, \ell_2, \ell_3)$ with $\ell_i \leq 3$ and $|\ell_1 - \ell_2| \leq \ell_3 \leq \ell_1 + \ell_2$, we (i) sample random input tensors of matching ranks, (ii) convert to the spherical basis via the orthonormal Cartesian-to-spherical change of basis of \citet{shao2025orthonormal}, (iii) apply the $\ell_1 \otimes \ell_2 \to \ell_3$ CG coupling using standard Wigner $3j$-symbols, (iv) convert back to the Cartesian basis, and (v) compare against the output of our constructive map. The maximum absolute error across all tested channels and $10^3$ random samples is below $10^{-11}$ in double precision; numerical logs are included with the code release.

\textbf{End-to-end equivariance test.} As a complementary end-to-end check, a smoke-test script applies a random $\mathrm{SO}(3)$ rotation $Q$ to every input position, rebuilds the radius graph, and verifies that (i) energies are unchanged ($|E(\mathrm{pos}) - E(Q\,\mathrm{pos})| < 10^{-3}$) and (ii) predicted forces rotate consistently ($\lVert F(Q\,\mathrm{pos}) - Q F(\mathrm{pos})\rVert_\infty < 10^{-3}$) across all ablation configurations. Equivariance is confirmed to expected numerical precision for every variant in Table~\ref{tab:ablation}.

\subsection{Edge embedding}\label{supp:edge_embed}

Pairwise distances $d_{ij} = \lVert\mathbf{r}_{ij}\rVert$ are expanded through a bank of Gaussian radial basis functions (RBFs) centred on evenly spaced reference distances between 0 and the cutoff $r_c$, followed by a smooth cosine envelope $\phi_c(d_{ij}) = \tfrac{1}{2}[\cos(\pi d_{ij}/r_c) + 1]$ that drives edge features smoothly to zero at the cutoff. Two small MLPs map the RBF expansion to channel-wise scalar weights for the grade-0 part and channel-wise weights for the grade-1 part of the edge multivector; the grade-1 direction $\hat{r}_{ij}$ is then multiplied by the channel weights and placed in the vector component. For CliffordSTF variants with $\mathrm{stf\_mode} \neq \text{none}$, a third MLP produces channel weights for the STF$_2$ edge feature, populated with $\mathrm{STF}_2(\hat{r}_{ij}, \hat{r}_{ij})$ (a canonical $L{=}2$ descriptor of the edge direction). When $\mathrm{stf\_mode}=\text{stf2+stf3}$, the STF$_3$ edge block is initialised to zero and is populated only via message passing within node features, since direction alone does not carry independent $L{=}3$ information.

\subsection{Augmented product at the many-body stage}\label{supp:aug_prod}

At the many-body stage of each interaction layer, an augmented product couples the Clifford and STF tracks through six bilinear terms. Given operands $a$ and $b$ each carrying Clifford, STF$_2$, and (optionally) STF$_3$ components, the augmented product computes:
\begin{itemize}
    \item The ordinary geometric product $a^{\mathrm{Cl}} b^{\mathrm{Cl}}$ via the $8\times8$ Cayley table;
    \item A new STF$_2$ from the grade-1 components via $\mathrm{STF}_2(a^{\mathrm{vec}}, b^{\mathrm{vec}})$;
    \item An updated STF$_2$ via the $L{=}2 \otimes L{=}1 \to L{=}2$ path, applied from either operand (symmetrised-then-detraced cross-contraction of $a^{\mathrm{STF}_2}$ with $b^{\mathrm{vec}}$, and vice versa);
    \item An $L{=}0$ contribution to grade-0 via $\langle a^{\mathrm{STF}_2}, b^{\mathrm{STF}_2}\rangle_F = a_{ij}^{\mathrm{STF}_2} b_{ij}^{\mathrm{STF}_2}$;
    \item A new STF$_3$ via $\mathrm{STF}_3(a^{\mathrm{STF}_2}, b^{\mathrm{vec}})$;
    \item An STF$_2$ update via $\mathrm{STF}_3 \otimes L{=}1 \to L{=}2$ contraction of $a^{\mathrm{STF}_3}$ against $b^{\mathrm{vec}}$.
\end{itemize}
Each term corresponds bijectively to a CG channel in one of the decompositions $\mathbf{1}\otimes\mathbf{1}$, $\mathbf{2}\otimes\mathbf{1}$, $\mathbf{2}\otimes\mathbf{2}$, or $\mathbf{3}\otimes\mathbf{1}$ enumerated in Supplementary~\ref{supp:cg_proof}. The cross-track gradient pathway created by these bilinears is what provides a usable training signal to the STF parameters, as demonstrated by the ablation (\S\ref{sec:results_scalar}).

\subsection{Equivariant normalisation and gating}\label{supp:normalization}

Each CliffordSTFLinear layer is block-diagonal across tracks: it applies independent linear maps of shape $C_{\mathrm{in}} \to C_{\mathrm{out}}$ over channels within the Clifford multivector, the STF$_2$ tensor, and the STF$_3$ tensor, without mixing them linearly. Cross-track information flow occurs only through the bilinear products described in \S\ref{sec:architecture}.

Track-wise equivariant normalisation rescales each track by a per-channel invariant norm. For the Clifford track we use the standard per-grade Clifford norm $\lVert h^{(k)}\rVert$ computed as the scalar part of $h^{(k)}\widetilde{h^{(k)}}$, where $\widetilde{\cdot}$ is the reverse \citep{doran2003geometric}; for STF$_2$ we use the Frobenius norm $\lVert S\rVert_F^2 = S_{ij}S_{ij}$ with $S_{zz} = -S_{xx}-S_{yy}$ substituted; for STF$_3$ we use the analogous Frobenius norm with multiplicity factors reflecting the symmetric tensor structure. All three norms are $\mathrm{SO}(3)$-invariant scalars.

Norm-gated activation applies SiLU to grade-0 features and a sigmoid-gated rescaling to each $L \geq 1$ track: for a feature $f$ at rank $\ell \geq 1$, we compute $g = \sigma(\mathrm{MLP}(\lVert f\rVert))$ with the sigmoid MLP acting on the invariant norm, and output $g \cdot f$. Since $g$ is a scalar and $f$ is equivariant, the product is equivariant. This is the norm-gating scheme of \citet{weiler20183d,schutt2021painn} adapted to the three tracks.

\subsection{Adaptive \texorpdfstring{$L$}{L}-track routing}\label{supp:routing}

Adaptive routing reduces effective angular bandwidth for atoms in symmetric environments. It is applied as a per-atom, per-track multiplicative gate $g_\ell \in [0, 1]$ between message-passing layers, rescaling the STF$_\ell$ track as $h^{\mathrm{STF}_\ell}_i \leftarrow g_\ell^{(i)} \cdot h^{\mathrm{STF}_\ell}_i$. The Clifford track is always active. Three modes are supported.
\begin{itemize}
    \item \textbf{None.} All atoms use all tracks; $g_\ell^{(i)} = 1$.
    \item \textbf{Static.} Gates are learned per atomic type via a small embedding table. Useful when chemically distinct atoms have systematically different angular-resolution needs.
    \item \textbf{Learned.} Gates are produced by a two-layer MLP with sigmoid output applied to four per-atom invariant descriptors: (i) coordination number $|\mathcal{N}(i)|$, (ii) angular variance $1 - \lVert\bar{r}_i\rVert$ where $\bar{r}_i = \tfrac{1}{|\mathcal{N}(i)|}\sum_{j\in\mathcal{N}(i)}\hat{r}_{ij}$, (iii) mean neighbour distance, (iv) grade-0 feature norm averaged over channels. All four descriptors are $\mathrm{SO}(3)$-invariant, so the gate is invariant and the gated product remains equivariant.
\end{itemize}

\subsection{Full ablation configurations}\label{supp:variants}

CliffordSTF is a single configurable architecture parameterised by five ablation flags: $\mathrm{stf\_mode} \in \{\text{none}, \text{stf2}, \text{stf2+stf3}\}$, $\mathrm{use\_hodge\_forces} \in \{\text{off}, \text{on}\}$, $\mathrm{use\_adaptive\_routing} \in \{\text{off}, \text{on}\}$, $\mathrm{routing\_mode} \in \{\text{none}, \text{static}, \text{learned}\}$, and $\mathrm{use\_cross\_track} \in \{\text{off}, \text{on}\}$. Table~\ref{tab:ablation} enumerates the eleven named variants used in the ablation study, alongside the vanilla $L{=}1$ Clifford baseline (which is recovered bit-for-bit by setting all five flags to their off/none state and is shown as the twelfth row of Table~\ref{tab:ablation} for comparison). Ten of the eleven variants are flag-parameterised; the eleventh, \emph{STF only output}, is a structural ablation outside the flag space: it removes the STF tracks from the message-passing scaffold entirely while retaining an STF-augmented force-readout head, isolating whether the directional gain comes from $L\!\geq\!2$ representations operating throughout message passing or from an $L\!\geq\!2$ readout alone.

% =========================================================
% Experiments-cited sections
% =========================================================

\subsection{Caveats and per-model exclusions for headline tables}\label{supp:results_caveats}

This subsection holds the per-model reasoning behind the caveat marks ($\ast$, $\dagger$, gray rows) used in Tables~\ref{tab:md17_fcos} and~\ref{tab:oc22}. The marks themselves are explained in one phrase each in the table captions; this subsection expands each justification.

\textbf{Strictly invariant models (SchNet, DimeNet++).} Both architectures are scalar-only and compute forces as the negative gradient of the predicted energy with respect to atomic positions, rather than from a direct-force head. The cosine similarity of gradient-derived forces depends on the energy--force loss balancing during training; under our uniform protocol this balance is not tuned per architecture, and gradient-derived directional fidelity is not a comparable diagnostic against models that predict forces directly. We render their rows in gray in Table~\ref{tab:md17_fcos} and exclude them from the per-architecture directional narrative in \S\ref{sec:results_md17}.

\textbf{FAENet.} FAENet predicts forces directly from a learned head but achieves equivariance via stochastic frame averaging rather than through architectural constraints on the message-passing operations themselves. Because the frames are sampled at training and inference time, per-sample force predictions are not strictly equivariant and per-sample directional fidelity is not strictly defined. We list FAENet in its natural $L\!\leq\!1$ cohort with an $\ast$ mark and exclude it from the per-architecture directional comparison.

\textbf{GotenNet.} GotenNet is an $L\!\geq\!2$ architecture ($\mathtt{lmax}{=}2$) with a direct-force head, and would naturally belong in the top block of Table~\ref{tab:md17_fcos}. Under our fixed-budget protocol it did not converge on rMD17. The published GotenNet recipe matches our learning rate ($10^{-4}$) but additionally relies on a $10{,}000$-step learning-rate warmup, an on-plateau learning-rate scheduler, and early-stopping patience $150$, none of which are part of our fixed-budget protocol (constant LR, patience $30$, no warm-up). We report the GotenNet result in its natural $L$-cohort with an $\ast$ mark for completeness and exclude it from the per-architecture directional narrative; the underperformance reflects protocol mismatch rather than architecture quality at recommended hyperparameters.

\textbf{EquiformerV2 absolute-energy prediction.} EquiformerV2 predicts absolute (non-referenced) molecular energies rather than energy differences relative to a per-element reference. Its rMD17 energy MAE therefore reflects the magnitude of the absolute binding energy ($\sim$247\,eV, consistent across all three $L_{\max}$ settings) rather than a prediction error, and is not directly comparable to the energy-referenced MAEs of the other rows. We mark this with $\dagger$ in Table~\ref{tab:md17_fcos} and in the corresponding panel of Supplementary~\ref{supp:md17_permol}; the force-cosine similarity column, which is independent of the energy-referencing convention, remains directly comparable across all rows.

\textbf{Fixed-budget non-convergence (MACE, ICTP).} Under our uniform $10^6$-parameter budget, MACE and ICTP at every reported $L_{\max}$ attain force MAE in the 3{,}000--6{,}000\,meV/\AA{} range on rMD17 and energy MAE in the 8--22\,eV range on OC22. Both architectures are designed to operate at substantially larger capacity (typically $5{-}20\times 10^6$ parameters under recommended hyperparameters), and the underperformance reported here reflects the fixed-budget protocol rather than architecture quality at native scale. We report the values honestly in the tables as part of the within-protocol comparison the paper makes (\S\ref{sec:protocol}); readers interested in MACE/ICTP at recommended hyperparameters should consult the original publications.

\subsection{Training protocol and parameter budgets}\label{supp:protocol}

\textbf{Parameter budgets and shared architectural hyperparameters.} All Clifford and CliffordSTF models evaluated in the main text use at most $1.06\times 10^6$ parameters. Channel counts are tuned from $C{=}80$ (base Clifford, 8D features) down to $C{=}60$ for the $D{=}13$ and $D{=}20$ STF variants to keep the parameter budget within the $10^6$ target. All other architectural hyperparameters are held fixed across variants: cutoff $r_c = 6.0$\,\AA, 50 Gaussian RBF bases, maximum 50 neighbours per atom, 5 message-passing layers, body order up to 4. Exact per-model parameter counts are listed in Table~\ref{tab:wallclock}.

\textbf{Shared optimisation settings.} All models are trained with Adam at learning rate $10^{-4}$, no weight decay, early stopping at patience 30, and no EMA or warm-up, as stated in \S\ref{sec:protocol}. Details are shared in Table \ref{tab:training-details}.

\begin{table}[t]
  \centering
  \small
  \caption{Per-benchmark training settings. All runs use Adam with base learning rate $1{\times}10^{-4}$, no weight decay, and early-stopping patience of $30$ epochs (\S4.1). OC20-S2EF and OC22-S2EF share the same energy+force loss formulation; the OC20-IS2RE and OC22-IS2RE rows are the deliberate exception (single scalar adsorption energy, L1 loss). $\lambda_F$ is the force-loss weight.}
  \label{tab:training-details}
  \resizebox{\textwidth}{!}{%
  \begin{tabular}{llrrlllr}
    \toprule
    Benchmark & Task & Epochs & Batch & LR schedule & Loss & $\lambda_F$ & Seeds \\
    \midrule
    QM9        & scalar --- 6 properties\textsuperscript{$\ast$}      & 250 & 128 & constant                                & MSE                                 & ---  & 3 \\
    Molecule3D & scalar --- HOMO--LUMO gap                            & 100 & 20  & StepLR (step $10$, $\gamma{=}0.8$)      & L1                                  & ---  & 3 \\
    rMD17       & energy + forces (10 molecules)                        & 250 & 32  & constant                                & MSE                                 & 1.0  & 3$^{\dagger}$ \\
    OC20       & IS2RE                                                & 50  & 16  & constant                                & L1                                  & ---  & 3 \\
    OC20       & S2EF                                                 & 50  & 16  & constant                                & per-atom MAE\,(E) $+$ $\ell_2$\,(F) & 3.0  & 3 \\
    OC22       & IS2RE                                                & 150 & 16  & constant                                & L1                                  & ---  & 3 \\
    OC22       & S2EF                                                 & 50  & 16  & constant                                & per-atom MAE\,(E) $+$ $\ell_2$\,(F) & 3.0  & 3 \\
    \bottomrule
  \end{tabular}}

  \vspace{0.4em}
  \begin{minipage}{\linewidth}
    \footnotesize
    $^{\ast}$QM9 properties: $\mu$, $\alpha$, $\varepsilon_\mathrm{HOMO}$, $\varepsilon_\mathrm{LUMO}$, $U_0$, $C_v$, trained as independent single-target models.$^{\dagger}$rMD17 Clifford + CliffordSTF ablation (Table~\ref{tab:ablation}): 12 architectural variants tuned to ${\sim}1$M parameters, evaluated on aspirin, benzene, ethanol, and salicylic acid with 6 seeds per (variant, molecule) obtained as 3 random seeds $\times$ 2 independent instances. Epochs, batch size, LR schedule, loss, force weight, and patience match the main rMD17 row.\emph{Preprocessing.} QM9 uses Z-score target normalisation and an 80/10/10 random split. rMD17 uses a 900/100/1000 train/validation/test random split; forces are predicted by a direct force head when the model exposes one (e.g.\ PaiNN, GotenNet) and are otherwise computed as the negative gradient of the energy with respect to atomic positions. Molecule3D uses a single random split over the ${\sim}780$k-molecule subset. OC20 and OC22 use a $6.0$~\AA{} neighbour cutoff; the S2EF rows additionally train and evaluate on free (non-fixed) atoms only, while the IS2RE rows do not apply this filter. OC22-IS2RE is trained for longer ($150$ epochs vs.\ $50$ for the other OC tasks) because the smaller scalar training signal converges more slowly. Each OC run uses official $200{,}000$ training split; the final evaluation reported in the paper uses the full validation and held-out splits.
  \end{minipage}
\end{table}

\subsection{Model architecture hyperparameters per benchmark}\label{supp:model_settings}

Tables~\ref{tab:supp_settings_md17}--\ref{tab:supp_settings_md17_ablation} list the architectural hyperparameters used per benchmark for every model in the main study. Cutoffs are in \AA{}; Boolean flags are abbreviated T/F. Training-time settings (optimiser, epochs, batch size, loss, seeds) are reported separately in Table~\ref{tab:training-details}; per-model parameter counts are in Table~\ref{tab:wallclock}.

\subsubsection{MD17}

Table~\ref{tab:supp_settings_md17} lists the per-model architecture hyperparameters used on rMD17. MACE, ICTP, and EquiformerV2 are run at three $L_{\max}$ values to trace angular resolution within each family at the matched ${\sim}10^6$-parameter budget, with EquiformerV2 reducing its sphere channel count from $40$ at $L_{\max}\!=\!2$ to $14$ at $L_{\max}\!=\!4$ to absorb the additional irreps. CliffordSTF runs at $C\!=\!60$ channels (versus $C\!=\!48$ for the base $L\!\leq\!1$ Clifford configuration on MD17) so that the STF$_2$ and STF$_3$ tracks fit within the same parameter envelope; cutoffs are $5$\,\AA{} for the equivariant transformers (TorchMD-Net, ViSNet) and NequIP, and $6$\,\AA{} elsewhere.

\begin{table}[h]
\centering
\caption{Model architecture hyperparameters used on \textbf{MD17}. All runs use Adam, no weight decay, batch size 32, learning rate 0.0001, 250 epochs, MSE loss, early-stopping patience 30. Cutoff is in \AA; Boolean flags are abbreviated T/F.}
\label{tab:supp_settings_md17}
\footnotesize
\setlength{\tabcolsep}{4pt}
\resizebox{\textwidth}{!}{%
\begin{tabular}{l c l}
\toprule
Model & Cutoff & Architecture \\
\midrule
SchNet & 6 & hidden=192, filters=192, layers=6, gauss=50 \\
PaiNN & 6 & hidden=144, layers=4, rbf=64 \\
DimeNet++ & 6 & hidden=128, blocks=4, radial=6, spherical=7, basis emb=8, int emb=64, out emb=128 \\
FAENet & 6 & hidden=320, filters=128, layers=4, FA=3D, FA-method=stochastic, MP-type=updownscale\_base \\
GotenNet & 6 & basis=96, layers=5 \\
TorchMD-Net & 5 & emb=104, layers=6, heads=8, rbf=56, rbf type=expnorm, arch=equivariant-transformer \\
ViSNet & 5 & hidden=96, layers=6, heads=8, rbf=32, $L_{\max}$=1, vertex=F \\
NequIP & 5 & features=44, layers=4, $L_{\max}$=1, bessels=8, parity=T \\
MACE L=1 & 6 & hidden=62x0e + 62x1o, layers=2, max ell=1, corr=3, bessel=8 \\
MACE L=2 & 6 & hidden=49x0e + 49x1o + 49x2e, layers=2, max ell=2, corr=3, bessel=8 \\
MACE L=3 & 6 & hidden=28x0e + 28x1o + 28x2e + 28x3o, layers=2, max ell=3, corr=3, bessel=8 \\
ICTP L=1 & 6 & hidden=10, layers=2, $L_{\max}$ hid=1, $L_{\max}$ edge=3, corr=3, basis=8, prod=10 \\
ICTP L=2 & 6 & hidden=8, layers=2, $L_{\max}$ hid=2, $L_{\max}$ edge=3, corr=3, basis=8, prod=8 \\
ICTP L=3 & 6 & hidden=7, layers=2, $L_{\max}$ hid=3, $L_{\max}$ edge=3, corr=3, basis=8, prod=7 \\
EquiformerV2 L=2 & 6 & sphere=40, layers=2, heads=4, $L_{\max}$=2, $M_{\max}$=2, attn $\alpha$=16, attn hid=20, attn val=8, ffn=64 \\
EquiformerV2 L=3 & 6 & sphere=24, layers=2, heads=4, $L_{\max}$=3, $M_{\max}$=2, attn $\alpha$=16, attn hid=20, attn val=8, ffn=64 \\
EquiformerV2 L=4 & 6 & sphere=14, layers=2, heads=4, $L_{\max}$=4, $M_{\max}$=2, attn $\alpha$=16, attn hid=20, attn val=8, ffn=64 \\
Clifford & 6 & channels=48, layers=5, rbf=50, hidden out=48 \\
CliffordSTF & 6 & channels=60, layers=5, rbf=50, STF mode=stf2+stf3, routing=learned, adapt rt=T, cross=T, Hodge F=T, GP rdo=F, self-int=F \\
\bottomrule
\end{tabular}%
}
\end{table}

\subsubsection{OC20 \& OC22 (IS2RE \& S2EF)}

Table~\ref{tab:supp_settings_oc} lists the architectures used for OC20 and OC22 on both IS2RE and S2EF; the per-model architectural hyperparameters are identical between the two benchmarks (same channel widths, layer counts, $L_{\max}$ choices, and Clifford-track configuration on every model), so the OC20 $\to$ OC22 comparison isolates dataset effects rather than architectural differences. The two benchmarks differ only in their per-model training schedules, captured in the caption. All models share a $6$\,\AA{} cutoff; MACE and ICTP are run at $L_{\max}\!\in\!\{1, 2\}$ (the $L\!=\!3$ configurations from rMD17 are dropped because they did not fit within the budget at OC graph density), and EquiformerV2 uses a single $L_{\max}\!=\!4$ configuration with a smaller sphere size (sphere=$18$, versus $14$--$40$ on rMD17) to keep parameter count near $10^6$. The base Clifford row is configured with $\mathrm{max\_grade}\!=\!1$ and the $\ell\!=\!2$ flag disabled, recovering the $L\!\leq\!1$ multivector substrate that the CliffordSTF $L\!\geq\!2$ extension is benchmarked against in the main results.

\begin{table}[h]
\centering
\caption{Model architecture hyperparameters used on \textbf{OC20 \& OC22 (IS2RE \& S2EF)}. The per-model architecture is identical between OC20 and OC22; only training schedules differ. \textbf{OC20:} all runs use Adam, no weight decay, batch size 16, learning rate 0.0001, 50 epochs, early-stopping patience 30; L1 loss on IS2RE for all models, MSE loss on S2EF for all models. \textbf{OC22:} all runs use Adam, no weight decay, batch size 16, learning rate 0.0001; on IS2RE, 200 epochs and patience 50; on S2EF, 50 epochs, MSE loss, patience 30 for all models. Cutoff is in \AA; Boolean flags are abbreviated T/F.}
\label{tab:supp_settings_oc}
\footnotesize
\setlength{\tabcolsep}{4pt}
\resizebox{\textwidth}{!}{%
\begin{tabular}{l c l}
\toprule
Model & Cutoff & Architecture \\
\midrule
SchNet & 6 & hidden=192, filters=192, layers=6, gauss=50 \\
PaiNN & 6 & hidden=144, layers=4, rbf=64 \\
DimeNet++ & 6 & hidden=128, blocks=4, radial=6, spherical=7, basis emb=8, int emb=64, out emb=128 \\
GotenNet & 6 & basis=96, layers=5 \\
TorchMD-Net & 6 & emb=104, layers=6, heads=8, rbf=56, rbf type=expnorm, arch=equivariant-transformer \\
NequIP & 6 & features=44, layers=4, $L_{\max}$=1, bessels=8, parity=T \\
MACE L=1 & 6 & hidden=62x0e + 62x1o, layers=2, max ell=1, corr=3, bessel=8 \\
MACE L=2 & 6 & hidden=49x0e + 49x1o + 49x2e, layers=2, max ell=2, corr=3, bessel=8 \\
ICTP L=1 & 6 & hidden=10, layers=2, $L_{\max}$ hid=1, $L_{\max}$ edge=3, corr=3, basis=8, prod=10 \\
ICTP L=2 & 6 & hidden=8, layers=2, $L_{\max}$ hid=2, $L_{\max}$ edge=3, corr=3, basis=8, prod=8 \\
EquiformerV2 & 6 & sphere=18, layers=2, heads=4, $L_{\max}$=4, $M_{\max}$=2, attn $\alpha$=16, attn hid=20, attn val=8, ffn=64 \\
Clifford & 6 & channels=80, layers=5, rbf=50, hidden out=80, max grade=1, $\ell{=}2$=F \\
CliffordSTF & 6 & channels=60, layers=5, rbf=50, STF mode=stf2+stf3, routing=learned, adapt rt=T, cross=T, Hodge F=T, GP rdo=F, self-int=F \\
\bottomrule
\end{tabular}%
}
\end{table}

\subsubsection{QM9}

Table~\ref{tab:supp_settings_qm9} lists the architectures used for the QM9 scalar-target benchmark. The QM9 model set is a subset of the rMD17 lineup: MACE and ICTP are not evaluated, EquiformerV2 is reported at a single $L_{\max}\!=\!4$ configuration with a larger spherical-channel count (sphere=$30$) than its catalysis runs, and the base $L\!\leq\!1$ Clifford row is omitted because the directional diagnostic that motivates its inclusion on rMD17 does not bear directly on scalar prediction. CliffordSTF uses the same $C\!=\!60$ channel count and the same STF$_2$+STF$_3$ configuration as on the other benchmarks.

\begin{table}[h]
\centering
\caption{Model architecture hyperparameters used on \textbf{QM9}. All runs use Adam, no weight decay, batch size 128, learning rate 0.0001, 250 epochs, MSE loss, early-stopping patience 30. Cutoff is in \AA; Boolean flags are abbreviated T/F.}
\label{tab:supp_settings_qm9}
\footnotesize
\setlength{\tabcolsep}{4pt}
\resizebox{\textwidth}{!}{%
\begin{tabular}{l c l}
\toprule
Model & Cutoff & Architecture \\
\midrule
SchNet & 6 & hidden=192, filters=192, layers=6, gauss=50 \\
PaiNN & 6 & hidden=144, layers=4, rbf=64 \\
DimeNet++ & 6 & hidden=128, blocks=4, radial=6, spherical=7, basis emb=8, int emb=64, out emb=128 \\
FAENet & 6 & hidden=320, filters=128, layers=4, FA=3D, FA-method=stochastic, MP-type=updownscale\_base \\
GotenNet & 6 & basis=96, layers=5 \\
TorchMD-Net & 5 & emb=104, layers=6, heads=8, rbf=56, rbf type=expnorm, arch=equivariant-transformer \\
ViSNet & 5 & hidden=96, layers=6, heads=8, rbf=32, $L_{\max}$=1, vertex=F \\
NequIP & 5 & features=44, layers=4, $L_{\max}$=1, bessels=8, parity=T \\
EquiformerV2 & 6 & sphere=30, layers=2, heads=4, $L_{\max}$=4, $M_{\max}$=2, attn $\alpha$=16, attn hid=20, attn val=8, ffn=64 \\
CliffordSTF & 6 & channels=60, layers=5, rbf=50, STF mode=stf2+stf3, routing=learned, adapt rt=T, cross=T, Hodge F=T, GP rdo=F, self-int=F \\
\bottomrule
\end{tabular}%
}
\end{table}

\subsubsection{Molecule3D}

Table~\ref{tab:supp_settings_molecule3d} lists the architectures used for the Molecule3D HOMO--LUMO gap benchmark. The model set is identical to QM9 with ViSNet additionally not reported; per-model architectural hyperparameters are otherwise unchanged from QM9 (same channel counts, layer counts, and $L_{\max}$ values per model), isolating the inter-benchmark difference to dataset scale and target rather than architecture.

\begin{table}[h]
\centering
\caption{Model architecture hyperparameters used on \textbf{Molecule3D}. All runs use Adam, no weight decay, batch size 20, learning rate 0.0001, 100 epochs, L1 loss, early-stopping patience 30. Cutoff is in \AA; Boolean flags are abbreviated T/F.}
\label{tab:supp_settings_molecule3d}
\footnotesize
\setlength{\tabcolsep}{4pt}
\resizebox{\textwidth}{!}{%
\begin{tabular}{l c l}
\toprule
Model & Cutoff & Architecture \\
\midrule
SchNet & 6 & hidden=192, filters=192, layers=6, gauss=50 \\
PaiNN & 6 & hidden=144, layers=4, rbf=64 \\
DimeNet++ & 6 & hidden=128, blocks=4, radial=6, spherical=7, basis emb=8, int emb=64, out emb=128 \\
FAENet & 6 & hidden=320, filters=128, layers=4, FA=3D, FA-method=stochastic, MP-type=updownscale\_base \\
GotenNet & 6 & basis=96, layers=5 \\
TorchMD-Net & 5 & emb=104, layers=6, heads=8, rbf=56, rbf type=expnorm, arch=equivariant-transformer \\
NequIP & 5 & features=44, layers=4, $L_{\max}$=1, bessels=8, parity=T \\
EquiformerV2 & 6 & sphere=30, layers=2, heads=4, $L_{\max}$=4, $M_{\max}$=2, attn $\alpha$=16, attn hid=20, attn val=8, ffn=64 \\
CliffordSTF & 6 & channels=60, layers=5, rbf=50, STF mode=stf2+stf3, routing=learned, adapt rt=T, cross=T, Hodge F=T, GP rdo=F, self-int=F \\
\bottomrule
\end{tabular}%
}
\end{table}

\subsubsection{MD17 ablation}

Table~\ref{tab:supp_settings_md17_ablation} lists the architectural hyperparameters of the twelve CliffordSTF ablation variants together with the vanilla $L\!\leq\!1$ Clifford baseline used in the controlled rMD17 ablation (Table~\ref{tab:ablation}). Channel counts are tuned per variant ($C\!\in\!\{60, 64, 76, 80\}$) so that every variant lands within $\pm 50$\% of the $10^6$-parameter target; the five ablation flags ($\mathrm{stf\_mode}$, $\mathrm{routing\_mode}$, $\mathrm{adapt\_rt}$, $\mathrm{cross}$, $\mathrm{Hodge\_F}$) toggle the architectural extensions documented in Supplementary~\ref{supp:variants}. Setting all flags to their off/none state recovers the base Clifford row bit-for-bit, providing the clean $L\!\leq\!1$ baseline against which every CliffordSTF extension is measured.

\begin{table}[h]
\centering
\caption{Model architecture hyperparameters used on \textbf{MD17 ablation}. All runs use Adam, no weight decay, batch size 32, learning rate 0.0001, 250 epochs, MSE loss, early-stopping patience 30. Cutoff is in \AA; Boolean flags are abbreviated T/F.}
\label{tab:supp_settings_md17_ablation}
\footnotesize
\setlength{\tabcolsep}{4pt}
\resizebox{\textwidth}{!}{%
\begin{tabular}{l c l}
\toprule
Model & Cutoff & Architecture \\
\midrule
STF (full) & 6 & channels=60, layers=5, rbf=50, STF mode=stf2+stf3, routing=learned, adapt rt=T, cross=T, Hodge F=T, GP rdo=F, self-int=F \\
STF$_2$ (static routing) & 6 & channels=64, layers=5, rbf=50, STF mode=stf2, routing=static, adapt rt=T, cross=T, Hodge F=T, GP rdo=F, self-int=F \\
STF baseline & 6 & channels=76, layers=5, rbf=50, STF mode=none, routing=none, adapt rt=F, cross=F, Hodge F=F, GP rdo=F, self-int=F \\
STF$_2$ (learned routing) & 6 & channels=64, layers=5, rbf=50, STF mode=stf2, routing=learned, adapt rt=T, cross=T, Hodge F=T, GP rdo=F, self-int=F \\
STF$_2$+STF$_3$ & 6 & channels=60, layers=5, rbf=50, STF mode=stf2+stf3, routing=none, adapt rt=F, cross=T, Hodge F=T, GP rdo=F, self-int=F \\
STF (full, no cross-track) & 6 & channels=60, layers=5, rbf=50, STF mode=stf2+stf3, routing=none, adapt rt=F, cross=F, Hodge F=T, GP rdo=F, self-int=F \\
STF (Hodge only) & 6 & channels=76, layers=5, rbf=50, STF mode=none, routing=none, adapt rt=F, cross=F, Hodge F=T, GP rdo=F, self-int=F \\
STF$_2$ (no Hodge) & 6 & channels=64, layers=5, rbf=50, STF mode=stf2, routing=none, adapt rt=F, cross=T, Hodge F=F, GP rdo=F, self-int=F \\
STF$_2$ (no cross-track) & 6 & channels=64, layers=5, rbf=50, STF mode=stf2, routing=none, adapt rt=F, cross=F, Hodge F=T, GP rdo=F, self-int=F \\
STF$_2$ & 6 & channels=64, layers=5, rbf=50, STF mode=stf2, routing=none, adapt rt=F, cross=T, Hodge F=T, GP rdo=F, self-int=F \\
STF (full, no routing) & 6 & channels=60, layers=5, rbf=50, STF mode=stf2+stf3, routing=none, adapt rt=F, cross=T, Hodge F=T, GP rdo=F, self-int=F \\
STF (output only) & 6 & channels=56, layers=5, rbf=50, STF mode=stf2+stf3, routing=none, adapt rt=F, cross=T, Hodge F=F, GP rdo=F, self-int=F \\
Clifford (vanilla) & 6 & channels=80, layers=5, rbf=50, hidden out=80, max grade=1, $\ell{=}2$=F \\
\bottomrule
\end{tabular}%
}
\end{table}

\subsection{rMD17 per-molecule force and energy MAE}\label{supp:md17_permol}

Table~\ref{tab:md17_permol_full} reports per-molecule force MAE (meV/\AA) and energy MAE (meV) for all models in the main-text rMD17 evaluation, in native units and averaged across seeds with standard deviations, complementing the force cosine similarity table in the main paper (Table~\ref{tab:md17_fcos}). The Clifford-to-CliffordSTF force cosine improvement observed in Table~\ref{tab:md17_fcos} occurs at comparable force and energy magnitude accuracy, which is the content of the "matched magnitude" claim in the main text. Reading panel (a), CliffordSTF's $17.6$\,meV/\AA{} aggregate force MAE sits between the converged $L\!\geq\!2$ tier (EquiformerV2 at $3.9$--$4.2$\,meV/\AA{} across $L_{\max}\!\in\!\{2,3,4\}$) and the base Clifford row at $20.8$\,meV/\AA{}, and improves on the base Clifford on every individual molecule. Panel (b) shows the energy-MAE picture parallels the force result: CliffordSTF and the base Clifford track each other within seed variance ($30.4$ vs.\ $34.1$\,meV aggregate), confirming that the STF extension does not trade scalar accuracy for directional accuracy.

\begin{table}[h]
\centering
\caption{\textbf{rMD17 per-molecule force and energy MAE} (mean $\pm$ std across 3 seeds, native units: meV/\AA{} for force MAE, meV for energy MAE). Values $\geq 1{,}000$ are rendered in k-notation (e.g., \texttt{12.8k $\pm$ 2.0k} $\equiv$ \texttt{12{,}800 $\pm$ 2{,}000}) to keep converged-model rows visually distinct from non-converged baselines. Top block: $L\!\geq\!2$ direct-force models. Middle block: $L\!\leq\!1$ direct-force models. Bottom block (gray): strictly invariant models (SchNet, DimeNet++) whose forces come from energy gradients. $\ast$ FAENet predicts forces directly but uses stochastic frame-averaged equivariance; GotenNet is $L\!\geq\!2$ ($\mathtt{lmax}{=}2$) but did not converge under our fixed-budget protocol---both are listed in their natural $L$-cohort for completeness and excluded from the per-architecture directional comparison in the main text. Under the fixed $10^6$-parameter budget, MACE and ICTP at every $L_{\max}$ attain force MAE in the 3{,}000--6{,}000\,meV/\AA{} range, reflecting non-convergence at this capacity rather than architecture quality at recommended hyperparameters. $\dagger$ EquiformerV2 predicts absolute (non-referenced) molecular energies; its energy MAE in panel (b) reflects absolute binding-energy magnitudes ($\sim$247\,eV consistent across all three $L_{\max}$ settings), not prediction error, and is not comparable to the energy-referenced baselines.}
\label{tab:md17_permol_full}
\setlength{\tabcolsep}{3pt}
\scriptsize
{(a) Force MAE (meV/\AA).}\\[2pt]
\resizebox{\textwidth}{!}{%
\begin{tabular}{lccccccccccc}
\toprule
Model & Asp & Azo & Bnz & Eth & Mal & Nap & Par & Sal & Tol & Ura & Agg \\
\midrule
EquiformerV2 $L{=}4$ & 6.172 $\pm$ 0.239 & 4.451 $\pm$ 0.274 & 1.805 $\pm$ 0.072 & 2.137 $\pm$ 0.022 & 3.193 $\pm$ 0.135 & 3.066 $\pm$ 0.058 & 5.653 $\pm$ 0.157 & 4.869 $\pm$ 0.755 & 3.298 $\pm$ 0.310 & 4.649 $\pm$ 0.259 & 3.929 $\pm$ 0.104 \\
EquiformerV2 $L{=}3$ & 6.290 $\pm$ 0.283 & 4.237 $\pm$ 0.131 & 1.917 $\pm$ 0.214 & 2.605 $\pm$ 0.220 & 4.246 $\pm$ 0.146 & 2.973 $\pm$ 0.066 & 5.308 $\pm$ 0.221 & 4.872 $\pm$ 0.220 & 3.316 $\pm$ 0.202 & 4.369 $\pm$ 0.048 & 4.013 $\pm$ 0.081 \\
EquiformerV2 $L{=}2$ & 6.281 $\pm$ 0.239 & 4.563 $\pm$ 0.329 & 1.830 $\pm$ 0.077 & 3.775 $\pm$ 0.299 & 4.569 $\pm$ 0.549 & 2.945 $\pm$ 0.094 & 5.292 $\pm$ 0.322 & 5.117 $\pm$ 0.178 & 3.611 $\pm$ 0.241 & 4.483 $\pm$ 0.180 & 4.247 $\pm$ 0.108 \\
CliffordSTF (ours) & 21.419 $\pm$ 0.453 & 17.944 $\pm$ 2.527 & 9.692 $\pm$ 0.675 & 15.302 $\pm$ 4.451 & 17.696 $\pm$ 3.218 & 16.657 $\pm$ 2.360 & 20.071 $\pm$ 1.390 & 19.820 $\pm$ 1.137 & 16.113 $\pm$ 2.414 & 20.930 $\pm$ 0.907 & 17.564 $\pm$ 0.671 \\
MACE $L{=}3$ & 6.1k $\pm$ 1.0k & 4.9k $\pm$ 1.1k & 1.0k $\pm$ 0.2k & 1.4k $\pm$ 0.2k & 2.6k $\pm$ 0.6k & 3.1k $\pm$ 0.4k & 4.9k $\pm$ 1.2k & 6.0k $\pm$ 1.8k & 2.6k $\pm$ 0.3k & 4.9k $\pm$ 0.2k & 3.7k $\pm$ 0.4k \\
MACE $L{=}2$ & 8.2k $\pm$ 0.9k & 5.3k $\pm$ 1.0k & 1.3k $\pm$ 0.2k & 2.1k $\pm$ 0.5k & 2.9k $\pm$ 0.1k & 5.2k $\pm$ 1.1k & 7.7k $\pm$ 1.3k & 6.4k $\pm$ 2.2k & 4.4k $\pm$ 0.9k & 4.9k $\pm$ 0.5k & 4.8k $\pm$ 0.2k \\
ICTP $L{=}3$ & 12.8k $\pm$ 2.0k & 6.0k $\pm$ 2.1k & 1.7k $\pm$ 0.8k & 1.9k $\pm$ 0.1k & 3.9k $\pm$ 0.6k & 5.6k $\pm$ 0.8k & 7.7k $\pm$ 0.9k & 7.8k $\pm$ 2.5k & 3.4k $\pm$ 0.3k & 3.7k $\pm$ 0.7k & 5.4k $\pm$ 0.6k \\
ICTP $L{=}2$ & 7.5k $\pm$ 1.5k & 6.9k $\pm$ 1.0k & 2.3k $\pm$ 0.2k & 1.9k $\pm$ 0.2k & 2.3k $\pm$ 1.1k & 5.6k $\pm$ 1.4k & 7.7k $\pm$ 2.8k & 8.2k $\pm$ 1.9k & 3.5k $\pm$ 1.0k & 7.1k $\pm$ 0.6k & 5.3k $\pm$ 0.2k \\
GotenNet$^{\ast}$ & 28.1k $\pm$ 24.0k & 4.3k $\pm$ 3.2k & 14.0k $\pm$ 22.1k & 2.2k $\pm$ 2.7k & 7.5k $\pm$ 12.3k & 3.2k $\pm$ 1.6k & 8.3k $\pm$ 10.0k & 3.9k $\pm$ 3.7k & 2.5k $\pm$ 3.3k & 14.3k $\pm$ 14.8k & 8.8k $\pm$ 2.7k \\
\midrule
Clifford (ours base) & 21.855 $\pm$ 0.002 & 22.262 $\pm$ 0.010 & 15.212 $\pm$ 0.009 & 19.996 $\pm$ 0.001 & 21.922 $\pm$ 0.004 & 21.345 $\pm$ 0.024 & 21.313 $\pm$ 0.004 & 21.470 $\pm$ 0.027 & 21.029 $\pm$ 0.004 & 22.074 $\pm$ 0.037 & 20.848 $\pm$ 0.003 \\
MACE $L{=}1$ & 10.7k $\pm$ 1.3k & 5.8k $\pm$ 0.7k & 2.6k $\pm$ 0.4k & 2.3k $\pm$ 0.3k & 4.5k $\pm$ 0.4k & 6.7k $\pm$ 0.3k & 6.4k $\pm$ 1.0k & 8.3k $\pm$ 1.4k & 5.5k $\pm$ 2.3k & 6.3k $\pm$ 1.2k & 5.9k $\pm$ 0.4k \\
ICTP $L{=}1$ & 10.4k $\pm$ 4.3k & 5.7k $\pm$ 0.7k & 2.2k $\pm$ 0.5k & 2.0k $\pm$ 0.4k & 4.0k $\pm$ 1.1k & 4.8k $\pm$ 1.0k & 8.4k $\pm$ 1.3k & 7.5k $\pm$ 1.6k & 3.1k $\pm$ 0.2k & 4.5k $\pm$ 3.2k & 5.3k $\pm$ 0.6k \\
PaiNN & 1.4k $\pm$ 0.2k & 1.1k $\pm$ 0.0k & 387.093 $\pm$ 20.028 & 392.650 $\pm$ 50.219 & 662.121 $\pm$ 190.588 & 804.449 $\pm$ 239.850 & 970.710 $\pm$ 382.994 & 1.0k $\pm$ 0.2k & 784.899 $\pm$ 178.245 & 1.1k $\pm$ 0.3k & 863.323 $\pm$ 7.657 \\
TorchMD-NET & 640.948 $\pm$ 78.214 & 347.636 $\pm$ 20.053 & 154.117 $\pm$ 7.309 & 154.106 $\pm$ 1.947 & 299.581 $\pm$ 97.919 & 219.310 $\pm$ 38.192 & 486.259 $\pm$ 81.638 & 386.224 $\pm$ 85.780 & 234.554 $\pm$ 31.349 & 278.833 $\pm$ 25.750 & 320.157 $\pm$ 8.706 \\
ViSNet & 1.7k $\pm$ 0.7k & 1.7k $\pm$ 0.7k & 539.034 $\pm$ 77.909 & 425.601 $\pm$ 309.210 & 825.760 $\pm$ 267.674 & 832.077 $\pm$ 276.489 & 1.0k $\pm$ 0.1k & 1.3k $\pm$ 0.1k & 573.367 $\pm$ 152.801 & 1.1k $\pm$ 0.2k & 996.860 $\pm$ 43.575 \\
NequIP $L{=}1$ & 6.1k $\pm$ 0.0k & 2.7k $\pm$ 0.1k & 1.2k $\pm$ 0.0k & 1.2k $\pm$ 0.0k & 2.4k $\pm$ 0.0k & 1.8k $\pm$ 0.0k & 4.4k $\pm$ 0.0k & 2.9k $\pm$ 0.0k & 1.1k $\pm$ 0.0k & 2.5k $\pm$ 0.0k & 2.6k $\pm$ 0.0k \\
FAENet$^{\ast}$ & 21.886 $\pm$ 0.030 & 22.288 $\pm$ 0.053 & 11.807 $\pm$ 2.931 & 19.980 $\pm$ 0.008 & 21.814 $\pm$ 0.179 & 21.548 $\pm$ 0.867 & 26.748 $\pm$ 9.464 & 21.679 $\pm$ 0.098 & 20.840 $\pm$ 0.301 & 21.553 $\pm$ 1.018 & 21.014 $\pm$ 0.749 \\
\midrule
\color{gray}{DimeNet++} & \color{gray}{1.2k $\pm$ 0.1k} & \color{gray}{991.164 $\pm$ 223.096} & \color{gray}{278.098 $\pm$ 49.866} & \color{gray}{389.387 $\pm$ 103.601} & \color{gray}{726.850 $\pm$ 80.483} & \color{gray}{608.244 $\pm$ 3.732} & \color{gray}{1.1k $\pm$ 0.1k} & \color{gray}{988.503 $\pm$ 107.070} & \color{gray}{519.458 $\pm$ 69.633} & \color{gray}{1.0k $\pm$ 0.1k} & \color{gray}{776.368 $\pm$ 11.972} \\
\color{gray}{SchNet} & \color{gray}{1.2k $\pm$ 0.1k} & \color{gray}{918.727 $\pm$ 33.117} & \color{gray}{196.534 $\pm$ 20.019} & \color{gray}{97.289 $\pm$ 13.538} & \color{gray}{226.752 $\pm$ 38.836} & \color{gray}{700.782 $\pm$ 58.526} & \color{gray}{1.4k $\pm$ 0.1k} & \color{gray}{1.2k $\pm$ 0.0k} & \color{gray}{627.706 $\pm$ 87.829} & \color{gray}{664.972 $\pm$ 35.163} & \color{gray}{720.170 $\pm$ 4.806} \\
\bottomrule
\end{tabular}%
}
 
\vspace{8pt}
\scriptsize
{(b) Energy MAE (meV).}\\[2pt]
\resizebox{\textwidth}{!}{%
\begin{tabular}{lccccccccccc}
\toprule
Model & Asp & Azo & Bnz & Eth & Mal & Nap & Par & Sal & Tol & Ura & Agg \\
\midrule
EquiformerV2 $L{=}4$$^{\dagger}$ & 405.9k $\pm$ 0.0k & 358.3k $\pm$ 0.0k & 145.3k $\pm$ 0.0k & 96.9k $\pm$ 0.0k & 167.2k $\pm$ 0.0k & 241.4k $\pm$ 0.0k & 322.5k $\pm$ 0.1k & 310.4k $\pm$ 0.0k & 169.8k $\pm$ 0.0k & 259.6k $\pm$ 0.0k & 247.7k $\pm$ 0.0k \\
EquiformerV2 $L{=}3$$^{\dagger}$ & 405.8k $\pm$ 0.0k & 358.2k $\pm$ 0.1k & 145.1k $\pm$ 0.0k & 96.9k $\pm$ 0.0k & 167.1k $\pm$ 0.0k & 241.2k $\pm$ 0.0k & 322.4k $\pm$ 0.0k & 310.3k $\pm$ 0.1k & 169.7k $\pm$ 0.1k & 259.5k $\pm$ 0.0k & 247.6k $\pm$ 0.0k \\
EquiformerV2 $L{=}2$$^{\dagger}$ & 405.5k $\pm$ 0.1k & 357.9k $\pm$ 0.0k & 145.0k $\pm$ 0.0k & 96.8k $\pm$ 0.0k & 167.0k $\pm$ 0.0k & 241.0k $\pm$ 0.0k & 322.1k $\pm$ 0.0k & 310.1k $\pm$ 0.1k & 169.5k $\pm$ 0.0k & 259.4k $\pm$ 0.0k & 247.4k $\pm$ 0.0k \\
CliffordSTF (ours) & 28.054 $\pm$ 9.040 & 46.565 $\pm$ 20.081 & 17.160 $\pm$ 2.661 & 6.677 $\pm$ 1.432 & 25.323 $\pm$ 5.299 & 22.087 $\pm$ 14.892 & 60.367 $\pm$ 30.139 & 25.819 $\pm$ 18.136 & 39.777 $\pm$ 20.171 & 31.876 $\pm$ 3.656 & 30.370 $\pm$ 6.565 \\
MACE $L{=}3$ & 2.6k $\pm$ 0.3k & 1.5k $\pm$ 0.5k & 188.788 $\pm$ 42.079 & 393.426 $\pm$ 84.371 & 775.686 $\pm$ 134.804 & 930.046 $\pm$ 54.572 & 2.4k $\pm$ 1.1k & 1.9k $\pm$ 0.9k & 957.564 $\pm$ 141.704 & 1.4k $\pm$ 0.1k & 1.3k $\pm$ 0.2k \\
MACE $L{=}2$ & 4.2k $\pm$ 1.2k & 1.9k $\pm$ 0.3k & 411.629 $\pm$ 325.937 & 653.666 $\pm$ 354.932 & 833.440 $\pm$ 22.310 & 1.9k $\pm$ 0.4k & 3.6k $\pm$ 0.8k & 2.3k $\pm$ 1.0k & 1.8k $\pm$ 0.3k & 1.5k $\pm$ 0.5k & 1.9k $\pm$ 0.3k \\
ICTP $L{=}3$ & 5.9k $\pm$ 0.3k & 2.6k $\pm$ 0.6k & 365.226 $\pm$ 27.122 & 515.550 $\pm$ 60.480 & 1.1k $\pm$ 0.3k & 1.9k $\pm$ 0.4k & 3.3k $\pm$ 0.6k & 2.7k $\pm$ 0.6k & 1.2k $\pm$ 0.0k & 1.1k $\pm$ 0.3k & 2.1k $\pm$ 0.2k \\
ICTP $L{=}2$ & 3.2k $\pm$ 0.5k & 3.1k $\pm$ 0.7k & 356.589 $\pm$ 14.078 & 516.129 $\pm$ 179.758 & 662.887 $\pm$ 233.720 & 2.0k $\pm$ 0.8k & 4.2k $\pm$ 2.0k & 3.1k $\pm$ 0.9k & 1.4k $\pm$ 0.8k & 2.2k $\pm$ 0.3k & 2.1k $\pm$ 0.1k \\
GotenNet$^{\ast}$ & 21.2k $\pm$ 17.1k & 3.8k $\pm$ 4.3k & 20.0k $\pm$ 33.9k & 1.1k $\pm$ 1.6k & 104.7k $\pm$ 79.1k & 2.4k $\pm$ 1.2k & 54.1k $\pm$ 90.2k & 14.5k $\pm$ 23.8k & 1.1k $\pm$ 1.3k & 37.4k $\pm$ 49.6k & 26.0k $\pm$ 7.2k \\
\midrule
Clifford (ours base) & 66.537 $\pm$ 31.283 & 56.820 $\pm$ 44.698 & 24.707 $\pm$ 17.343 & 14.675 $\pm$ 6.332 & 37.370 $\pm$ 25.576 & 14.780 $\pm$ 5.935 & 23.000 $\pm$ 7.962 & 35.300 $\pm$ 29.797 & 39.009 $\pm$ 31.720 & 29.097 $\pm$ 17.892 & 34.130 $\pm$ 2.563 \\
MACE $L{=}1$ & 4.9k $\pm$ 0.6k & 2.2k $\pm$ 0.6k & 446.266 $\pm$ 78.046 & 748.730 $\pm$ 294.846 & 1.4k $\pm$ 0.3k & 2.5k $\pm$ 0.1k & 2.9k $\pm$ 0.6k & 3.2k $\pm$ 0.1k & 2.4k $\pm$ 1.1k & 1.9k $\pm$ 0.3k & 2.3k $\pm$ 0.1k \\
ICTP $L{=}1$ & 5.1k $\pm$ 2.6k & 2.2k $\pm$ 0.5k & 709.468 $\pm$ 201.768 & 646.268 $\pm$ 219.460 & 1.2k $\pm$ 0.4k & 1.9k $\pm$ 0.2k & 4.7k $\pm$ 0.8k & 3.0k $\pm$ 0.9k & 1.5k $\pm$ 0.1k & 1.4k $\pm$ 0.7k & 2.2k $\pm$ 0.4k \\
PaiNN & 4.5k $\pm$ 1.3k & 3.1k $\pm$ 3.5k & 654.920 $\pm$ 484.394 & 572.296 $\pm$ 497.010 & 1.4k $\pm$ 0.8k & 1.4k $\pm$ 1.1k & 2.6k $\pm$ 3.3k & 697.934 $\pm$ 132.654 & 1.4k $\pm$ 0.7k & 2.0k $\pm$ 1.5k & 1.8k $\pm$ 0.5k \\
TorchMD-NET & 1.4k $\pm$ 0.7k & 496.497 $\pm$ 269.580 & 180.477 $\pm$ 170.513 & 187.650 $\pm$ 148.968 & 486.742 $\pm$ 231.389 & 451.681 $\pm$ 462.015 & 1.6k $\pm$ 0.5k & 624.017 $\pm$ 754.032 & 263.880 $\pm$ 284.126 & 940.992 $\pm$ 872.265 & 661.312 $\pm$ 130.339 \\
ViSNet & 2.7k $\pm$ 1.5k & 6.5k $\pm$ 2.7k & 413.191 $\pm$ 358.609 & 1.5k $\pm$ 1.6k & 348.909 $\pm$ 95.459 & 2.2k $\pm$ 1.0k & 1.7k $\pm$ 0.3k & 4.1k $\pm$ 2.9k & 922.044 $\pm$ 640.715 & 2.9k $\pm$ 0.8k & 2.3k $\pm$ 0.2k \\
NequIP $L{=}1$ & 3.7k $\pm$ 0.1k & 1.0k $\pm$ 0.1k & 197.342 $\pm$ 39.519 & 453.727 $\pm$ 39.250 & 916.506 $\pm$ 14.226 & 542.585 $\pm$ 12.952 & 3.1k $\pm$ 0.1k & 1.2k $\pm$ 0.1k & 411.931 $\pm$ 64.980 & 755.870 $\pm$ 109.032 & 1.2k $\pm$ 0.0k \\
FAENet$^{\ast}$ & 33.876 $\pm$ 15.568 & 26.864 $\pm$ 26.312 & 4.012 $\pm$ 1.777 & 3.480 $\pm$ 0.059 & 3.351 $\pm$ 0.161 & 47.213 $\pm$ 38.921 & 87.265 $\pm$ 123.215 & 32.828 $\pm$ 30.261 & 8.166 $\pm$ 3.405 & 5.399 $\pm$ 2.175 & 25.245 $\pm$ 12.952 \\
\midrule
\color{gray}{DimeNet++} & \color{gray}{1.8k $\pm$ 1.1k} & \color{gray}{3.7k $\pm$ 2.4k} & \color{gray}{216.364 $\pm$ 250.668} & \color{gray}{611.701 $\pm$ 618.476} & \color{gray}{1.0k $\pm$ 0.3k} & \color{gray}{956.485 $\pm$ 380.148} & \color{gray}{1.1k $\pm$ 0.7k} & \color{gray}{1.6k $\pm$ 1.4k} & \color{gray}{724.374 $\pm$ 424.645} & \color{gray}{1.2k $\pm$ 1.4k} & \color{gray}{1.3k $\pm$ 0.4k} \\
\color{gray}{SchNet} & \color{gray}{1.2k $\pm$ 0.8k} & \color{gray}{1.3k $\pm$ 1.0k} & \color{gray}{604.539 $\pm$ 447.764} & \color{gray}{403.595 $\pm$ 239.889} & \color{gray}{588.952 $\pm$ 533.312} & \color{gray}{1.7k $\pm$ 0.7k} & \color{gray}{1.2k $\pm$ 0.4k} & \color{gray}{2.1k $\pm$ 0.9k} & \color{gray}{441.554 $\pm$ 134.053} & \color{gray}{981.058 $\pm$ 333.775} & \color{gray}{1.1k $\pm$ 0.1k} \\
\bottomrule
\end{tabular}%
}
\end{table}

\subsection{OC20 and OC22 full results}\label{supp:oc_full}

This subsection reports the full OC20 and OC22 results (Table~\ref{tab:oc_full}), mean $\pm$ std across seeds in native units. Panel~(a) gives OC20 at the 200k scale: IS2RE MAE for all four splits ($\mathrm{val}_{\text{id}}$, $\mathrm{val}_{\text{ood-ads}}$, $\mathrm{val}_{\text{ood-cat}}$, $\mathrm{val}_{\text{ood-both}}$) together with S2EF ($e_\mathrm{MAE}$, $f_\mathrm{MAE}$, $f_\mathrm{cos}$) for $\mathrm{val}_{\text{id}}$ and $\mathrm{val}_{\text{ood-both}}$; the two remaining OOD splits (ood-ads, ood-cat) behave similarly to ood-both on S2EF and are omitted for compactness. Panel~(b) reproduces the OC22 results from the main text (Table~\ref{tab:oc22}) alongside OC20 for side-by-side reading. Reading across the two panels, CliffordSTF and the base Clifford row attain the strongest in-distribution S2EF energy MAE on both benchmarks (panel (a) OC20: $2.115$/$1.889$\,eV vs.\ $2.454$\,eV for EquiformerV2; panel (b) OC22: $2.866$/$2.844$\,eV vs.\ $3.369$\,eV); CliffordSTF additionally is the best $L\!\geq\!2$ architecture on OC22 IS2RE in-distribution ($4.089$\,eV vs.\ next-best $L\!\geq\!2$ GotenNet at $4.318$\,eV and EquiformerV2 at $L_{\max}\!=\!4$ at $4.905$\,eV) and on OC22 IS2RE OOD EwT ($0.312$ vs $0.234$ for EquiformerV2). EquiformerV2 retains the strongest S2EF force cosine similarity on OC20 in our protocol; on OC22 the two are within seed variance.

\begin{table}[h]
\centering
\caption{\textbf{OC20 and OC22 full results} (mean $\pm$ std across 3 seeds; native units: eV for IS2RE MAE and S2EF energy MAE, eV/\AA{} for S2EF force MAE, dimensionless for force cosine similarity). (a) OC20 at the 200k scale, with IS2RE across all four splits and S2EF for $\mathrm{val}_{\text{id}}$ and $\mathrm{val}_{\text{ood-both}}$. (b) OC22 at full scale, with IS2RE MAE for $\mathrm{val}_{\text{id}}$ and $\mathrm{val}_{\text{ood}}$, IS2RE EwT for $\mathrm{val}_{\text{ood}}$, and S2EF for $\mathrm{val}_{\text{id}}$ and $\mathrm{val}_{\text{ood}}$. Our models are in the top two rows of each panel. MACE and ICTP at $L_{\max}\!\in\!\{1,2\}$ attain energy MAE values 1--3$\times$ those of converged baselines under the fixed $10^6$-parameter budget, mirroring the non-convergence pattern documented on rMD17 (Table~\ref{tab:md17_permol_full}).}
\label{tab:oc_full}
\setlength{\tabcolsep}{3pt}
\scriptsize
{(a) OC20}\\[2pt]
\resizebox{\textwidth}{!}{%
\begin{tabular}{lcccccccccc}
\toprule
& \multicolumn{4}{c}{IS2RE MAE (eV)} & \multicolumn{3}{c}{S2EF $\mathrm{val}_{\text{id}}$} & \multicolumn{3}{c}{S2EF $\mathrm{val}_{\text{ood-both}}$} \\
\cmidrule(lr){2-5}\cmidrule(lr){6-8}\cmidrule(lr){9-11}
Model & $\mathrm{val}_{\text{id}}$ & $\mathrm{val}_{\text{ood-ads}}$ & $\mathrm{val}_{\text{ood-cat}}$ & $\mathrm{val}_{\text{ood-both}}$ & $e_\mathrm{MAE}$ (eV) & $f_\mathrm{MAE}$ (eV/\AA) & $f_\mathrm{cos}\uparrow$ & $e_\mathrm{MAE}$ (eV) & $f_\mathrm{MAE}$ (eV/\AA) & $f_\mathrm{cos}\uparrow$ \\
\midrule
CliffordSTF (ours) & 0.703 $\pm$ 0.008 & 0.814 $\pm$ 0.024 & 0.707 $\pm$ 0.003 & 0.753 $\pm$ 0.029 & 2.115 $\pm$ 0.059 & 0.060 $\pm$ 0.000 & 0.138 $\pm$ 0.003 & 3.486 $\pm$ 0.160 & 0.077 $\pm$ 0.002 & 0.125 $\pm$ 0.006 \\
Clifford (ours base) & 0.765 $\pm$ 0.025 & 0.907 $\pm$ 0.031 & 0.761 $\pm$ 0.022 & 0.849 $\pm$ 0.025 & 1.889 $\pm$ 0.014 & 0.064 $\pm$ 0.000 & 0.115 $\pm$ 0.003 & 3.510 $\pm$ 0.004 & 0.081 $\pm$ 0.000 & 0.095 $\pm$ 0.004 \\
EquiformerV2 & 0.801 $\pm$ 0.152 & 0.845 $\pm$ 0.081 & 0.776 $\pm$ 0.149 & 0.785 $\pm$ 0.078 & 2.454 $\pm$ 0.013 & 0.060 $\pm$ 0.000 & 0.170 $\pm$ 0.002 & 3.602 $\pm$ 0.126 & 0.074 $\pm$ 0.000 & 0.162 $\pm$ 0.005 \\
GotenNet & 0.681 $\pm$ 0.006 & 0.749 $\pm$ 0.018 & 0.661 $\pm$ 0.002 & 0.682 $\pm$ 0.011 & 3.650 $\pm$ 0.276 & 0.064 $\pm$ 0.003 & 0.167 $\pm$ 0.006 & 4.850 $\pm$ 0.165 & 0.082 $\pm$ 0.001 & 0.144 $\pm$ 0.007 \\
TorchMD-NET & 0.726 $\pm$ 0.022 & 0.826 $\pm$ 0.046 & 0.714 $\pm$ 0.030 & 0.762 $\pm$ 0.039 & 3.715 $\pm$ 0.006 & 0.064 $\pm$ 0.000 & 0.162 $\pm$ 0.003 & 4.931 $\pm$ 0.010 & 0.081 $\pm$ 0.000 & 0.140 $\pm$ 0.002 \\
PaiNN & 0.720 $\pm$ 0.005 & 0.900 $\pm$ 0.088 & 0.706 $\pm$ 0.008 & 0.821 $\pm$ 0.070 & 4.198 $\pm$ 0.002 & 0.066 $\pm$ 0.000 & 0.154 $\pm$ 0.000 & 5.339 $\pm$ 0.009 & 0.083 $\pm$ 0.000 & 0.132 $\pm$ 0.001 \\
DimeNet++ & 0.699 $\pm$ 0.013 & 0.843 $\pm$ 0.024 & 0.691 $\pm$ 0.014 & 0.771 $\pm$ 0.031 & 4.098 $\pm$ 0.023 & 0.063 $\pm$ 0.000 & 0.170 $\pm$ 0.000 & 5.373 $\pm$ 0.053 & 0.080 $\pm$ 0.001 & 0.153 $\pm$ 0.001 \\
SchNet & 0.766 $\pm$ 0.007 & 0.834 $\pm$ 0.047 & 0.755 $\pm$ 0.001 & 0.779 $\pm$ 0.038 & 4.790 $\pm$ 0.067 & 0.070 $\pm$ 0.001 & 0.128 $\pm$ 0.006 & 5.879 $\pm$ 0.079 & 0.085 $\pm$ 0.001 & 0.115 $\pm$ 0.004 \\
ICTP $L{=}2$ & 0.881 $\pm$ 0.028 & 1.108 $\pm$ 0.027 & 0.830 $\pm$ 0.012 & 0.972 $\pm$ 0.001 & 6.434 $\pm$ 0.146 & 0.077 $\pm$ 0.001 & 0.100 $\pm$ 0.006 & 7.608 $\pm$ 0.235 & 0.092 $\pm$ 0.001 & 0.090 $\pm$ 0.005 \\
ICTP $L{=}1$ & 0.881 $\pm$ 0.020 & 0.996 $\pm$ 0.015 & 0.842 $\pm$ 0.002 & 0.904 $\pm$ 0.033 & 6.032 $\pm$ 0.149 & 0.079 $\pm$ 0.003 & 0.098 $\pm$ 0.002 & 7.212 $\pm$ 0.127 & 0.094 $\pm$ 0.003 & 0.089 $\pm$ 0.004 \\
MACE $L{=}2$ & 0.936 $\pm$ 0.134 & 0.977 $\pm$ 0.174 & 0.900 $\pm$ 0.126 & 0.881 $\pm$ 0.132 & 6.268 $\pm$ 0.002 & 0.075 $\pm$ 0.001 & 0.113 $\pm$ 0.002 & 7.367 $\pm$ 0.020 & 0.091 $\pm$ 0.001 & 0.102 $\pm$ 0.001 \\
MACE $L{=}1$ & 0.938 $\pm$ 0.153 & 1.052 $\pm$ 0.016 & 0.910 $\pm$ 0.130 & 0.974 $\pm$ 0.015 & 5.708 $\pm$ 0.042 & 0.077 $\pm$ 0.001 & 0.101 $\pm$ 0.002 & 6.979 $\pm$ 0.108 & 0.092 $\pm$ 0.001 & 0.092 $\pm$ 0.002 \\
NequIP & 0.931 $\pm$ 0.048 & 0.971 $\pm$ 0.083 & 0.899 $\pm$ 0.053 & 0.893 $\pm$ 0.076 & 5.802 $\pm$ 0.065 & 0.077 $\pm$ 0.000 & 0.110 $\pm$ 0.001 & 7.112 $\pm$ 0.009 & 0.091 $\pm$ 0.000 & 0.100 $\pm$ 0.000 \\
\bottomrule
\end{tabular}%
}
 
\vspace{8pt}
{(b) OC22}\\[2pt]
\resizebox{\textwidth}{!}{%
\begin{tabular}{lccccccccc}
\toprule
& \multicolumn{3}{c}{IS2RE} & \multicolumn{3}{c}{S2EF $\mathrm{val}_{\text{id}}$} & \multicolumn{3}{c}{S2EF $\mathrm{val}_{\text{ood}}$} \\
\cmidrule(lr){2-4}\cmidrule(lr){5-7}\cmidrule(lr){8-10}
Model & $\mathrm{MAE}_{\text{id}}$ (eV) & $\mathrm{MAE}_{\text{ood}}$ (eV) & $\mathrm{EwT}_{\text{ood}}\uparrow$ & $e_\mathrm{MAE}$ (eV) & $f_\mathrm{MAE}$ (eV/\AA) & $f_\mathrm{cos}\uparrow$ & $e_\mathrm{MAE}$ (eV) & $f_\mathrm{MAE}$ (eV/\AA) & $f_\mathrm{cos}\uparrow$ \\
\midrule
CliffordSTF (ours) & 4.089 $\pm$ 0.261 & 6.463 $\pm$ 0.310 & 0.312 $\pm$ 0.136 & 2.866 $\pm$ 0.093 & 0.073 $\pm$ 0.000 & 0.050 $\pm$ 0.004 & 5.083 $\pm$ 0.166 & 0.070 $\pm$ 0.000 & 0.050 $\pm$ 0.003 \\
Clifford (ours base) & 4.422 $\pm$ 0.862 & 6.726 $\pm$ 0.656 & 0.264 $\pm$ 0.055 & 2.844 $\pm$ 0.100 & 0.072 $\pm$ 0.000 & 0.059 $\pm$ 0.002 & 5.375 $\pm$ 0.157 & 0.071 $\pm$ 0.001 & 0.058 $\pm$ 0.002 \\
EquiformerV2 & 4.905 $\pm$ 0.566 & 6.619 $\pm$ 0.445 & 0.234 $\pm$ 0.025 & 3.369 $\pm$ 0.066 & 0.072 $\pm$ 0.000 & 0.067 $\pm$ 0.001 & 5.262 $\pm$ 0.076 & 0.070 $\pm$ 0.000 & 0.062 $\pm$ 0.002 \\
PaiNN & 3.516 $\pm$ 0.292 & 6.712 $\pm$ 0.140 & 0.234 $\pm$ 0.178 & 5.306 $\pm$ 0.030 & 0.076 $\pm$ 0.000 & 0.048 $\pm$ 0.001 & 7.452 $\pm$ 0.254 & 0.078 $\pm$ 0.000 & 0.036 $\pm$ 0.001 \\
GotenNet & 4.318 $\pm$ 0.300 & 5.772 $\pm$ 0.532 & 0.144 $\pm$ 0.000 & 4.666 $\pm$ 0.310 & 0.073 $\pm$ 0.000 & 0.052 $\pm$ 0.001 & 7.200 $\pm$ 0.121 & 0.072 $\pm$ 0.001 & 0.044 $\pm$ 0.001 \\
TorchMD-NET & 5.397 $\pm$ 0.000 & 8.312 $\pm$ 0.000 & 0.396 $\pm$ 0.000 & 4.278 $\pm$ 0.000 & 0.075 $\pm$ 0.000 & 0.051 $\pm$ 0.000 & 6.232 $\pm$ 0.000 & 0.074 $\pm$ 0.000 & 0.043 $\pm$ 0.000 \\
DimeNet++ & 6.320 $\pm$ 4.141 & 9.897 $\pm$ 3.002 & 0.108 $\pm$ 0.102 & 5.288 $\pm$ 0.030 & 0.076 $\pm$ 0.000 & 0.054 $\pm$ 0.002 & 8.717 $\pm$ 0.075 & 0.077 $\pm$ 0.001 & 0.038 $\pm$ 0.002 \\
SchNet & 6.402 $\pm$ 1.688 & 9.424 $\pm$ 0.938 & 0.270 $\pm$ 0.076 & 7.690 $\pm$ 1.508 & 0.077 $\pm$ 0.002 & 0.032 $\pm$ 0.007 & 12.016 $\pm$ 1.938 & 0.078 $\pm$ 0.002 & 0.024 $\pm$ 0.003 \\
MACE $L{=}2$ & 8.416 $\pm$ 0.702 & 22.569 $\pm$ 7.674 & 0.144 $\pm$ 0.051 & 10.017 $\pm$ 0.395 & 0.079 $\pm$ 0.000 & 0.022 $\pm$ 0.004 & 15.415 $\pm$ 0.808 & 0.076 $\pm$ 0.000 & 0.022 $\pm$ 0.005 \\
ICTP $L{=}2$ & 8.641 $\pm$ 1.744 & 17.756 $\pm$ 4.145 & 0.108 $\pm$ 0.051 & 12.447 $\pm$ 0.354 & 0.081 $\pm$ 0.001 & 0.015 $\pm$ 0.002 & 19.218 $\pm$ 0.232 & 0.079 $\pm$ 0.001 & 0.017 $\pm$ 0.002 \\
ICTP $L{=}1$ & 10.963 $\pm$ 1.612 & 20.810 $\pm$ 0.868 & 0.198 $\pm$ 0.076 & 11.483 $\pm$ 0.881 & 0.080 $\pm$ 0.001 & 0.019 $\pm$ 0.001 & 17.263 $\pm$ 0.655 & 0.186 $\pm$ 0.154 & 0.019 $\pm$ 0.000 \\
MACE $L{=}1$ & 12.339 $\pm$ 4.540 & 15.800 $\pm$ 2.430 & 0.180 $\pm$ 0.102 & 10.011 $\pm$ 0.018 & 0.078 $\pm$ 0.000 & 0.024 $\pm$ 0.000 & 15.737 $\pm$ 0.438 & 0.076 $\pm$ 0.000 & 0.023 $\pm$ 0.001 \\
NequIP & 12.575 $\pm$ 4.889 & 14.865 $\pm$ 3.676 & 0.036 $\pm$ 0.000 & 9.868 $\pm$ 0.148 & 0.079 $\pm$ 0.000 & 0.028 $\pm$ 0.000 & 15.809 $\pm$ 0.071 & 0.076 $\pm$ 0.000 & 0.027 $\pm$ 0.001 \\
\bottomrule
\end{tabular}%
}
\end{table}
 

\subsection{QM9 and Molecule3D full results}\label{supp:qm9_full}

Per-target MAE on QM9 and Molecule3D HOMO--LUMO gap, reported in Table~\ref{tab:qm9_mol3d_full} with native units per column and rows ordered by QM9 aggregate. The QM9 aggregate column averages the six per-target MAEs across mixed units, providing a single scalar per model that is comparable within the column; CliffordSTF places fifth of ten models at $2.152$\,units, sitting alongside the vector-equivariant baselines (PaiNN $2.229$, TorchMD-NET $2.232$) and ahead of the next-best $L\!\geq\!2$ row (EquiformerV2 $3.327$). On the Molecule3D HOMO--LUMO gap, FAENet's $0.178$\,eV is an order of magnitude ahead of the rest of the field; the remaining models cluster between $1.107$ and $1.495$\,eV with CliffordSTF at $1.232$\,eV, and we treat Molecule3D as out-of-scope for the directional thesis (\S\ref{sec:limitations}).

\begin{table}[h]
\centering
\caption{\textbf{QM9 per-target MAE and Molecule3D HOMO--LUMO gap MAE} (mean $\pm$ std across seeds; native units). QM9 targets are reported in their native units: heat capacity $C_v$ in cal/(mol$\cdot$K), atomisation energy $U_0$ in eV, polarisability $\alpha$ in Bohr$^3$, HOMO/LUMO in eV, dipole moment $\mu$ in Debye. The QM9 aggregate column is the mean of the six per-target MAEs across mixed units and serves as a single comparable scalar for between-model comparison. Molecule3D reports HOMO--LUMO gap MAE in eV on the large-scale dataset. Rows ordered by QM9 aggregate (lower is better).}
\label{tab:qm9_mol3d_full}
\setlength{\tabcolsep}{3pt}
\resizebox{\textwidth}{!}{%
\begin{tabular}{lcccccccc}
\toprule
& \multicolumn{7}{c}{QM9} & Molecule3D \\
\cmidrule(lr){2-8}\cmidrule(lr){9-9}
Model & $C_v$ & $U_0$ & $\alpha$ & HOMO & LUMO & $\mu$ & Aggregate & HOMO--LUMO (eV) \\
\midrule
GotenNet & 0.041 $\pm$ 0.008 & 8.045 $\pm$ 1.812 & 0.088 $\pm$ 0.007 & 0.033 $\pm$ 0.001 & 0.029 $\pm$ 0.001 & 0.024 $\pm$ 0.003 & 1.377 $\pm$ 0.304 & 1.117 $\pm$ 0.010 \\
FAENet & 0.045 $\pm$ 0.006 & 9.704 $\pm$ 5.215 & 0.114 $\pm$ 0.006 & 0.042 $\pm$ 0.001 & 0.036 $\pm$ 0.003 & 0.061 $\pm$ 0.006 & 1.667 $\pm$ 0.866 & 0.178 $\pm$ 0.002 \\
DimeNet++ & 0.063 $\pm$ 0.014 & 10.843 $\pm$ 2.719 & 0.122 $\pm$ 0.025 & 0.033 $\pm$ 0.001 & 0.028 $\pm$ 0.001 & 0.052 $\pm$ 0.003 & 1.857 $\pm$ 0.449 & 1.258 $\pm$ 0.020 \\
ViSNet & 0.058 $\pm$ 0.007 & 11.742 $\pm$ 6.513 & 0.223 $\pm$ 0.072 & 0.034 $\pm$ 0.001 & 0.031 $\pm$ 0.000 & 0.062 $\pm$ 0.044 & 2.025 $\pm$ 1.087 & -- \\
CliffordSTF (ours) & 0.055 $\pm$ 0.018 & 12.588 $\pm$ 9.200 & 0.143 $\pm$ 0.023 & 0.044 $\pm$ 0.001 & 0.041 $\pm$ 0.002 & 0.041 $\pm$ 0.003 & 2.152 $\pm$ 1.535 & 1.232 $\pm$ 0.005 \\
PaiNN & 0.042 $\pm$ 0.006 & 13.128 $\pm$ 3.356 & 0.110 $\pm$ 0.013 & 0.035 $\pm$ 0.000 & 0.031 $\pm$ 0.001 & 0.027 $\pm$ 0.005 & 2.229 $\pm$ 0.560 & 1.211 $\pm$ 0.003 \\
TorchMD-NET & 0.053 $\pm$ 0.014 & 13.153 $\pm$ 6.127 & 0.101 $\pm$ 0.008 & 0.033 $\pm$ 0.000 & 0.028 $\pm$ 0.001 & 0.024 $\pm$ 0.002 & 2.232 $\pm$ 1.024 & 1.149 $\pm$ 0.004 \\
EquiformerV2 & 0.075 $\pm$ 0.010 & 19.622 $\pm$ 8.347 & 0.171 $\pm$ 0.037 & 0.032 $\pm$ 0.001 & 0.031 $\pm$ 0.003 & 0.031 $\pm$ 0.002 & 3.327 $\pm$ 1.397 & 1.107 $\pm$ 0.049 \\
SchNet & 0.091 $\pm$ 0.018 & 23.424 $\pm$ 4.389 & 0.172 $\pm$ 0.038 & 0.044 $\pm$ 0.001 & 0.045 $\pm$ 0.001 & 0.071 $\pm$ 0.007 & 3.975 $\pm$ 0.732 & 1.410 $\pm$ 0.008 \\
NequIP & 0.171 $\pm$ 0.003 & 35.516 $\pm$ 0.991 & 0.515 $\pm$ 0.034 & 0.115 $\pm$ 0.001 & 0.142 $\pm$ 0.004 & 0.186 $\pm$ 0.003 & 6.108 $\pm$ 0.160 & 1.495 $\pm$ 0.000 \\
\bottomrule
\end{tabular}%
}
\end{table}


\subsection{Wall-clock measurements at matched parameter budget}\label{supp:wallclock}

Table~\ref{tab:wallclock} reports training and inference wall-clock on rMD17 (aspirin, batch 32, single GPU) for every model in the main-text study, measured under a strict timing protocol (10 warmup steps, 20 timed steps, mean $\pm$ std reported on a single GPU; absolute times are not directly comparable to figures reported on other hardware). At inference, the base Clifford model runs at $0.128$\,s/step, ranking fifth across the nineteen-model study and outpacing every $L\!\geq\!2$ or Cartesian-tensor baseline in the table: it is faster than EquiformerV2 $L{=}4$ ($0.142$\,s) by roughly $10\%$, faster than every MACE configuration ($L{=}1$: $0.160$; $L{=}2$: $0.237$; $L{=}3$: $0.261$), faster than NequIP ($0.221$), ViSNet ($0.235$), TorchMD-NET ($0.186$), DimeNet++ ($0.326$), and every ICTP variant ($0.302$--$0.364$). The four models with faster inference (PaiNN, GotenNet, SchNet, FAENet) are scalar-invariant or $L{=}1$ vector architectures with substantially less representational capacity. CliffordSTF inference at $0.352$\,s/step is competitive with the slower spherical-harmonic baselines, faster than ICTP $L{=}2$/$L{=}3$ and within $\sim\!8\%$ of DimeNet++. Since deployed interatomic potentials run inference far more often than they are trained, this is the practitioner-relevant operating regime.

Training cost is also more competitive than previously reported. The base Clifford model trains at $3.11$\,s/step---within $\sim\!25\%$ of EquiformerV2 $L{=}4$ ($2.48$\,s), tied with NequIP ($3.06$\,s), and faster than ICTP at every $L$ ($1.45\times$, $1.73\times$, and $2.00\times$ faster than ICTP $L{=}1/2/3$ respectively), faster than DimeNet++ ($1.88\times$), and faster than MACE $L{=}3$. CliffordSTF training at $7.77$\,s/step is comparable to capacity-scaled $L\!\geq\!2$ baselines, sitting within $25$--$73\%$ of ICTP $L{=}1/2/3$ and DimeNet++ ($1.25\times$ ICTP $L{=}3$, $1.33\times$ DimeNet++, $1.44\times$ ICTP $L{=}2$, $1.73\times$ ICTP $L{=}1$); the $2$--$3\times$ overhead is concentrated against the fastest spherical-harmonic baselines (MACE $L{=}1/2$, EquiformerV2 $L{=}2$--$4$). Both Clifford-family rows share the optimized geometric-product dispatch described in Supplementary~\ref{supp:optim}; the optimization narrowed the CliffordSTF/EquiformerV2 $L{=}4$ training ratio from a previously reported $\sim\!4.3\times$ on different hardware to $\sim\!3.1\times$ here, and the corresponding base Clifford ratio from $1.85\times$ to $1.25\times$.

The forward/backward decomposition pinpoints where the residual CliffordSTF training overhead lives. CliffordSTF spends $2.59$\,s in forward and $5.10$\,s in backward (backward is $\sim\!2\times$ forward), whereas EquiformerV2 $L{=}4$ has a near-symmetric profile ($1.09$\,s forward, $1.32$\,s backward, $1.21\times$); MACE $L{=}2$ is essentially symmetric ($1.22$\,s\,/\,$1.20$\,s). The autograd graph through the dual-track scaffold---not the geometric-product dispatch itself, which the optimization in Supplementary~\ref{supp:optim} addressed---dominates the residual training cost. The trade-off between training cost and accuracy per parameter (where CliffordSTF is competitive with strong tensor-field baselines on OC22) is discussed in \S\ref{sec:limitations}.

\begin{table}[h]
\centering
\caption{\textbf{Wall-clock at matched parameter budget}, rMD17 aspirin, batch 32, single GPU. Time columns are mean $\pm$ standard deviation over 20 timed steps following 10 warmup steps; samples-per-second columns are derived from the means. Lower is faster for the time columns; higher is faster for the samples-per-second columns. All models fall within $\pm 50$\% of the $10^6$-parameter target. The Clifford and CliffordSTF (ours) entries reflect the optimized PyTorch dispatch described in Supplementary~\ref{supp:optim}; bit-for-bit preservation of model outputs is verified by O(3)-equivariance tests.}
\label{tab:wallclock}
\small
\begin{tabular}{lrrrrr}
\toprule
Model & Params & Train step (s) & Train (s/s) & Infer step (s) & Infer (s/s) \\
\midrule
PaiNN              & 1{,}043{,}856 & 0.79 $\pm$ 0.05 & 1{,}145 & 0.028 $\pm$ 0.003 & 3{,}575 \\
SchNet             &   984{,}769 & 1.16 $\pm$ 0.02 &   773 & 0.069 $\pm$ 0.011 & 1{,}456 \\
FAENet             & 1{,}027{,}745 & 1.69 $\pm$ 0.03 &   533 & 0.111 $\pm$ 0.004 &   900 \\
GotenNet           & 1{,}057{,}729 & 1.81 $\pm$ 0.00 &   497 & 0.064 $\pm$ 0.005 & 1{,}571 \\
MACE $L{=}1$       & 1{,}006{,}014 & 1.81 $\pm$ 0.29 &   497 & 0.160 $\pm$ 0.008 &   624 \\
EquiformerV2 $L{=}2$ &   891{,}560 & 2.45 $\pm$ 0.01 &   367 & 0.131 $\pm$ 0.001 &   764 \\
EquiformerV2 $L{=}3$ & 1{,}041{,}304 & 2.45 $\pm$ 0.01 &   367 & 0.135 $\pm$ 0.003 &   741 \\
EquiformerV2 $L{=}4$ & 1{,}043{,}222 & 2.48 $\pm$ 0.00 &   363 & 0.142 $\pm$ 0.008 &   702 \\
MACE $L{=}2$       & 1{,}011{,}683 & 2.53 $\pm$ 0.09 &   355 & 0.237 $\pm$ 0.002 &   422 \\
TorchMD-NET        &   968{,}242 & 2.58 $\pm$ 0.02 &   349 & 0.186 $\pm$ 0.029 &   538 \\
NequIP             & 1{,}002{,}436 & 3.06 $\pm$ 0.04 &   295 & 0.221 $\pm$ 0.007 &   452 \\
Clifford (ours)    & 1{,}031{,}926 & 3.11 $\pm$ 0.01 &   289 & 0.128 $\pm$ 0.000 &   781 \\
ViSNet             &   957{,}606 & 3.21 $\pm$ 0.07 &   280 & 0.235 $\pm$ 0.001 &   425 \\
MACE $L{=}3$       &   980{,}092 & 3.60 $\pm$ 0.16 &   250 & 0.261 $\pm$ 0.009 &   384 \\
ICTP $L{=}1$       & 1{,}039{,}554 & 4.50 $\pm$ 0.11 &   200 & 0.302 $\pm$ 0.052 &   331 \\
ICTP $L{=}2$       &   997{,}368 & 5.38 $\pm$ 0.10 &   167 & 0.364 $\pm$ 0.007 &   275 \\
DimeNet++          & 1{,}064{,}070 & 5.84 $\pm$ 0.11 &   154 & 0.326 $\pm$ 0.002 &   307 \\
ICTP $L{=}3$       &   982{,}752 & 6.23 $\pm$ 0.02 &   144 & 0.356 $\pm$ 0.006 &   281 \\
CliffordSTF (ours) & 1{,}056{,}534 & 7.77 $\pm$ 0.11 &   116 & 0.352 $\pm$ 0.001 &   284 \\
\bottomrule
\end{tabular}
\end{table}

\subsection{Implementation optimizations}\label{supp:optim}

The Clifford-family rows in Table~\ref{tab:wallclock} reflect a PyTorch-level refactor of the geometric-product dispatch, with no custom CUDA or Triton kernels written. The four largest changes are: (i) the Cl(3,0) geometric product is computed as a single \texttt{einsum} against a precomputed Cayley structure tensor, replacing 33 hand-unrolled scalar fused-multiply-adds that incurred one CUDA launch each in eager mode; (ii) per-layer \texttt{batch.max().item()} synchronization points were eliminated by threading the graph count through the forward pass as a keyword argument, allowing the GPU to run asynchronously across the full forward; (iii) layer-independent geometric invariants in the adaptive-routing module (coordination number, angular variance, mean edge distance) are precomputed once per forward rather than per layer; and (iv) validation-metric accumulation was moved on-device, with a single host transfer per epoch. All four changes preserve model outputs bit-for-bit, verified by O(3)-equivariance tests at fp32 precision.

Two alternative optimizations were investigated and rolled back. Engaging \texttt{torch.compile} required wrapping \texttt{torch\_scatter.scatter\_softmax} in a \texttt{@torch.compiler.disable} helper (Dynamo cannot FX-trace the custom op); even with the wrapper, Inductor recompiled subgraphs on every edge-count change under the dynamic-shape configuration required by \texttt{radius\_graph}, producing exponential recompile thrashing rather than forward progress. A second attempt rewrote the STF$_2$ and STF$_3$ products as outer-product einsums with symmetrize-then-extract: the forward pass became $\sim\!1\%$ faster, but the backward pass slowed by $\sim\!4\%$ because autograd retained the larger $3\times 3\times 3$ intermediate rather than the chain of small scalar tensors produced by the unrolled form. Both attempts left dormant scaffolding (a \texttt{use\_compile} configuration flag and the \texttt{@torch.compiler.disable} marker) for future revisits when upstream blockers are resolved.

Approximately $40$--$60\%$ additional speedup over the current measurements appears tractable through full \texttt{torch.compile} engagement (pending an upstream fix for \texttt{torch\_scatter} custom-op tracing or migration to \texttt{torch.scatter\_reduce\_}), CUDA Graphs for static-shape batches, and fused Triton kernels for the STF$_2$/STF$_3$ products and the scatter--softmax--multiply pattern. We additionally fixed an unrelated shape-transpose bug in our PaiNN force-head wrapper that was introduced during this optimization pass; the bug caused the model to crash rather than return incorrect outputs, so no published PaiNN result in this paper or its supplementary tables was produced from the buggy state.